% =============================================================
% Chapitre 4 : Conception de l'architecture Zero Trust
% Source : ../../CH4_Conception_de_l_architecture.md
% =============================================================
% -- Resserrement typographique local au chapitre 4 (valeurs restaurées en fin de fichier).
\setlength{\textfloatsep}{10pt plus 2pt minus 3pt}
\setlength{\floatsep}{8pt plus 2pt minus 2pt}
\setlength{\intextsep}{8pt plus 2pt minus 2pt}
\setlength{\abovecaptionskip}{4pt}
\setlength{\belowcaptionskip}{2pt}
% -- Placement des flottants : le chapitre porte onze figures et neuf tableaux ;
% les valeurs par defaut de LaTeX reservent une page entiere des qu'un flottant
% depasse la moitie de la justification, ce qui laissait ici du blanc autour de
% cinq figures. Valeurs restaurees en fin de fichier.
\renewcommand{\topfraction}{0.92}
\renewcommand{\bottomfraction}{0.75}
\renewcommand{\textfraction}{0.08}
\renewcommand{\floatpagefraction}{0.88}
\setcounter{topnumber}{3}
\setcounter{totalnumber}{4}

\chapter{Conception de l'architecture}
\label{ch:conception}
\soustitrechapitre{Vues d'architecture, plan d'identité et registre des décisions}

Ce chapitre présente l'architecture qui réalise les vingt exigences du \cref{ch:besoins-menaces}, chaque
affirmation étant adossée à une vue qui la rend vérifiable : une architecture qui ne documente que
ce qu'elle contient laisse penser que ce qui n'y figure pas a été oublié.

\section{Principes directeurs de conception}

Quatre principes ont guidé l'ensemble des décisions ; ils rendent chaque choix contestable.

\textbf{PR1 --- tout composant doit justifier son existence.} Un composant ajouté est à configurer,
mettre à jour, surveiller et payer : exploité par une seule personne, il dégrade la sécurité réelle
même s'il améliore la sécurité théorique. Règle appliquée : \textbf{un composant n'est retenu que
s'il couvre une menace du \cref{ch:besoins-menaces} qu'aucun composant existant ne couvre déjà} --- sept ont été
écartés (\cref{tab:composants-ecartes}).

\textbf{PR2 --- l'identité est le plan de contrôle, pas une couche.} Le modèle périmétrique segmente
par le réseau ; le modèle retenu segmente par l'identité, chaque charge de travail ayant la sienne.
La représenter comme une couche latérale reviendrait à admettre des chemins où elle n'est pas
consultée, soit la confiance implicite même (\cref{subsec:plans-transversaux}).

\textbf{PR3 --- un moteur de détection ne modifie jamais les preuves qu'il analyse.} Traduction de
la valeur métier VM5 et de l'exigence EX11, étendue au tableau de bord ; seule exception explicite,
l'enregistrement du verdict humain, donnée nouvelle et non modification.

\textbf{PR4 --- rien n'existe dans l'infrastructure qui n'existe d'abord dans le code.} Réponse à
la carence C2 et au point bloquant B5 de l'audit du \cref{ch:cadre-existant}, pour trois effets :
reproductibilité, contrôle avant existence, détection de dérive. Les commandes interactives restent
réservées à l'exploitation ponctuelle.

\section{Vues d'architecture}
\label{sec:vues-archi}

Cinq vues statiques décrivent le socle, le réseau étant donné dans ses deux états, avant et après
la refonte. Aucune ne peut afficher un composant qui n'est pas déclaré en code.

\subsection{Vue de contexte}

La \cref{fig:contexte} fixe la frontière du système et nomme ce dont le socle dépend sans le
maîtriser : six systèmes tiers, dont trois imports de données qui portent le risque de chaîne
d'approvisionnement --- dépendance ou image de base compromise, acteur A10 du \cref{ch:besoins-menaces}. Ces trois
imports sont différés, hors du chemin de la requête : une indisponibilité y dégrade
l'enrichissement ou la priorisation, jamais le service rendu. Le sixième, la zone de noms de
domaine tenue chez le bureau d'enregistrement, figure en dépendance externe et non en composant de
la couche de périmètre --- écart É10 : elle est consultée \emph{avant} tout contact avec le socle,
de sorte qu'une dérive y rend le socle injoignable sans qu'aucun contrôle interne ne s'en aperçoive.

\begin{figure}[htbp]
\centering
\begin{tikzpicture}[
  sys/.style={draw, very thick, rounded corners=3pt, fill=black!12, text width=4.6cm,
    align=center, font=\footnotesize, inner sep=5pt, minimum height=1.9cm},
  act/.style={draw, thick, rounded corners=2pt, fill=white, text width=2.8cm,
    align=center, font=\footnotesize, inner sep=3pt, minimum height=1.1cm},
  extr/.style={act, text width=3.1cm},
  sync/.style={fluxdonnee},
  asyn/.style={fluxevenement},
  etq/.style={font=\scriptsize, fill=white, inner sep=1pt, align=center}
]
\node[sys] (socle) at (7.4,3.4)
  {\textbf{Socle MENAL}\\ hébergement, supervision et chaîne de livraison\\
   {\scriptsize\color{black!55} deux projets cloud}};
\node[act]  (users) at (1.6,6.1) {Utilisateurs de l'application hébergée};
\node[act]  (soc)   at (1.6,3.4) {Analyste et administrateur};
\node[act]  (dev)   at (1.6,0.7) {Élève-ingénieur, exploitant unique};
\node[extr] (dns)   at (1.6,8.4) {Bureau d'enregistrement du domaine (zone de noms)};
\node[extr] (git)   at (13.2,6.1) {Dépôt de code et chaîne de livraison};
\node[extr] (mitre) at (13.2,3.4) {Référentiel de techniques d'attaque};
\node[extr] (cve)   at (13.2,0.7) {Catalogues de vulnérabilités exploitées};
\node[extr] (model) at (7.4,6.5)  {Dépôt public de modèles};
\node[extr] (totp)  at (7.4,0.3)  {Application d'authentification};
\draw[sync] (users.north) -- (dns.south) node[etq, midway, right=1pt] {résolution de noms};
\draw[sync] (users.east) -| (socle.north west) node[etq, pos=0.28] {HTTPS};
\draw[sync] (soc.east)   -- (socle.west);
\draw[sync] (dev.east)   -| (socle.south west) node[etq, pos=0.28] {exploitation};
\draw[sync] (git.west)   -| (socle.north east) node[etq, pos=0.28] {identité fédérée};
\draw[asyn] (mitre.west) -- (socle.east) node[etq, midway] {catalogue};
\draw[asyn] (cve.west)   -| (socle.south east) node[etq, pos=0.28] {gravité réelle};
\draw[sync] (totp)  -- (socle.south) node[etq, midway, right] {second facteur};
\draw[asyn] (model) -- (socle.north) node[etq, midway, right] {poids figés au \emph{build}};
\node[legendecadre, anchor=north west, font=\scriptsize] at (-0.3,-0.5)
  {\legendetrait{fluxdonnee}{échange synchrone, chemin de la requête}\quad
   \legendetrait{fluxevenement}{import différé, hors du chemin de la requête}};
\end{tikzpicture}
\caption[Vue de contexte : le socle, ses acteurs et les systèmes extérieurs]{Vue de contexte :
le socle, ses acteurs et les six systèmes extérieurs dont il dépend --- trois imports différés
portent le risque de chaîne d'approvisionnement, la zone de noms échappe seule à la description en
code (É10)}
\label{fig:contexte}
\end{figure}

\subsection{Modèle en couches et justification des deux plans transversaux}
\label{subsec:plans-transversaux}

L'architecture s'organise en sept niveaux numérotés (\cref{fig:modele-couches}), dont deux plans
transversaux plutôt que des couches ordinaires, suivant l'approche recommandée pour les
architectures cloud sécurisées~\cite{gcp_architecture_framework}. La conception initiale
comportait deux couches de même numéro et omettait l'enrichissement (L6) : la numérotation a été
révisée, un modèle en couches n'ayant de valeur que si documentation, code et tests désignent la
même chose.

\begin{figure}[htbp]
\centering
\begin{tikzpicture}[
  couche/.style={draw, very thick, rounded corners=1pt, fill=black!4, align=center,
    font=\footnotesize, text width=8.4cm, inner sep=3pt, minimum height=0.95cm},
  demi/.style={couche, text width=3.9cm},
  plan/.style={draw, very thick, dashed, rounded corners=1pt, fill=black!10, align=center,
    font=\footnotesize},
  fleche/.style={fluxdonnee},
  col/.style={fluxevenement},
  idl/.style={fluxcontrole}
]
\node[couche] (l1) at (4.9,7.98) {\textbf{L1 --- Périmètre d'entrée}\\ répartiteur HTTPS, filtrage applicatif, limitation de débit\\{\scriptsize\color{black!55} Cloud Load Balancing · Cloud Armor}};
\node[demi] (l3a) at (2.55,6.20) {\textbf{L3 --- Application hébergée}\\{\scriptsize\color{black!55} Cloud Run}};
\node[demi] (l3b) at (7.25,6.20) {\textbf{L3 --- API de supervision et tableau de bord}\\{\scriptsize\color{black!55} Cloud Run}};
\node[couche] (l4) at (4.9,4.6) {\textbf{L4 --- Réseau}\\ refus par défaut, accès privé, sortie contrôlée\\{\scriptsize\color{black!55} VPC · sortie réseau directe · règles de pare-feu}};
\node[demi] (l5a) at (2.55,3.2) {\textbf{L5 --- Base de données applicative}\\{\scriptsize\color{black!55} Cloud SQL for PostgreSQL}};
\node[demi] (l5b) at (7.25,3.2) {\textbf{L5 --- Entrepôt de supervision}\\{\scriptsize\color{black!55} BigQuery}};
\node[couche] (l6) at (4.9,1.8) {\textbf{L6 --- Enrichissement sémantique} (service d'encodage, sans sortie réseau)\\{\scriptsize\color{black!55} Cloud Run (interne)}};
\node[plan, text width=8.4cm, inner sep=3pt, minimum height=0.95cm] (l7) at (4.9,0.4)
  {\textbf{L7 --- OBSERVABILITÉ} (plan transversal) : collecte depuis toutes les couches\\{\scriptsize\color{black!55} Cloud Logging · Cloud Monitoring}};
\node[plan, minimum width=1.75cm, minimum height=9.20cm] (l2) at (10.6,4.28)
  {\rotatebox{90}{\parbox{7.6cm}{\centering\textbf{L2 --- IDENTITÉ ET SÉCURITÉ} (plan transversal) : vérifiée à chaque saut\\{\footnotesize\color{menaltrait} Cloud IAM · Secret Manager · Cloud KMS · Workload Identity Federation}}}};
\draw[fleche] (l1) -- (l3a);
\draw[fleche] (l1) -- (l3b);
\draw[fleche] (l3a) -- (l4);
\draw[fleche] (l3b) -- (l4);
\draw[fleche] (l4) -- (l5a);
\draw[fleche] (l4) -- (l5b);
\draw[fleche] (l4) -- (l6);
\draw[col] (-0.4,7.98) -- (-0.4,0.4);
\foreach \n in {l1,l3a,l4,l5a,l6}{\draw[col] (\n.west) -- (\n.west -| -0.4,0);}
\draw[col,-{Latex[length=1.8mm]}] (-0.4,0.4) -- (l7.west);
\draw[col,-{Latex[length=1.8mm]}] (l7.east) -- ++(0.35,0) |- (l5b.south east);
\node[font=\scriptsize, fill=white, inner sep=1pt] at (-0.4,5.3) {F4};
\foreach \n in {l1,l3b,l4,l5b,l6}{\draw[idl] (\n.east -| l2.west) -- (\n.east);}
\node[legendecadre, anchor=north west, font=\scriptsize] at (-0.65,-0.30)
  {\legendetrait{fluxdonnee}{traversée d'un saut}\quad
   \legendetrait{fluxevenement}{collecte des journaux (F4)}\quad
   \legendetrait{fluxcontrole}{vérification d'identité (L2)}};
\end{tikzpicture}
\caption[Modèle en couches : cinq couches empilées et deux plans transversaux]{Modèle en couches :
cinq couches empilées (L1, L3, L4, L5, L6) et deux plans transversaux (L2, L7), dont la bordure en
tirets marque la portée. La résolution de noms ne figure pas en L1 : elle est une dépendance
externe (\cref{fig:contexte}).}
\label{fig:modele-couches}
\end{figure}

\textbf{L2 est un plan} parce que l'identité est vérifiée à \emph{chaque} saut --- du périmètre
aux charges de travail, de celles-ci aux données, de la chaîne de livraison au
registre~\cite{nist800207} --- portée dont un bloc latéral ne rendrait pas compte ; \textbf{L7 en
est un} parce que la collecte part de \emph{toutes} les couches pour alimenter l'entrepôt de
supervision (L5), sans source de données autrement.

\subsection{Vue de déploiement et chemins réseau effectifs}

La \cref{fig:deploiement} situe les mêmes éléments comme composants réels reliés par des liens
réseau, dont le \emph{style} porte l'information la plus contre-intuitive du socle : tous les
chemins internes ne traversent pas le réseau privé, l'accès à l'entrepôt empruntant le chemin
managé du fournisseur, autorisé par l'identité seule et non par une règle réseau (É11). Le tableau
de bord, lui, ne rejoint pas l'API par un lien direct : son appel ressort par le répartiteur et
subit un \textbf{second passage du filtrage applicatif}, ce qui explique le comportement observé de
la limitation de débit sur les chemins d'authentification (\cref{subsec:incidents}).

\begin{figure}[htbp]
\centering
\begin{tikzpicture}[
  acteur/.style={draw, thick, rounded corners=3pt, fill=white, align=center,
    font=\footnotesize, text width=2.3cm, inner sep=3pt},
  noeud/.style={draw, thick, rounded corners=1pt, fill=black!4, align=center,
    font=\footnotesize, text width=2.3cm, inner sep=3pt},
  large/.style={noeud, text width=8.4cm},
  pub/.style={fluxdonnee},
  priv/.style={fluxdonnee, line width=1.3pt},
  man/.style={fluxcontrole},
  jour/.style={fluxevenement},
  fx/.style={font=\scriptsize\bfseries, fill=white, inner sep=1pt},
  etq/.style={font=\scriptsize, fill=white, inner sep=1pt}
]
\node[acteur] (users)   at (1.0,9.7) {Utilisateurs};
\node[acteur] (analyst) at (4.9,9.7) {Analyste / Admin};
\node[acteur] (repo)    at (13.0,9.7) {Dépôt de code};
\node[large] (perim) at (4.9,8.0) {\textbf{Périmètre (L1)} --- répartiteur HTTPS, filtrage applicatif, limitation de débit\\{\scriptsize\color{black!55} Cloud Load Balancing · Cloud Armor}};
\node[noeud, text width=3.4cm] (ci) at (13.0,8.0) {\textbf{Chaîne de livraison} fédérée, sans clé\\{\scriptsize\color{black!55} GitHub Actions · Workload Identity Federation}};
\node[noeud] (app)  at (1.0,6.2) {Application hébergée\\{\scriptsize\color{black!55} Cloud Run}};
\node[noeud] (api)  at (4.9,6.2) {API de supervision\\{\scriptsize\color{black!55} Cloud Run}};
\node[noeud] (dash) at (8.8,6.2) {Tableau de bord\\{\scriptsize\color{black!55} Cloud Run}};
\node[noeud, text width=2.6cm] (reg) at (13.0,6.2) {Registre d'images\\{\scriptsize\color{black!55} Artifact Registry}};
\node[large] (net) at (4.9,4.4) {\textbf{Réseau privé (L4)} --- refus par défaut, accès privé, sortie réseau directe\\{\scriptsize\color{black!55} VPC · sortie réseau directe · règles de pare-feu}};
\node[noeud] (db)    at (1.0,2.6) {Base de données\\{\scriptsize\color{black!55} Cloud SQL}};
\node[noeud, text width=2.6cm] (dw) at (5.4,2.6) {Entrepôt de supervision (SIEM)\\{\scriptsize\color{black!55} BigQuery}};
\node[noeud, text width=2.6cm] (sched) at (9.4,2.6) {Tâche planifiée\\d'enrichissement\\{\scriptsize\color{black!55} Cloud Scheduler · Cloud Run Job}};
\node[noeud, text width=2.6cm] (enc) at (13.0,2.6) {Service d'encodage (interne)\\{\scriptsize\color{black!55} Cloud Run (interne)}};
\node[large] (col) at (4.9,0.8) {\textbf{Collecte journalisée (L7)} --- quatre exports, filtrés à la source pour trois d'entre eux\\{\scriptsize\color{black!55} Cloud Logging}};
\draw[pub] (users) -- (users |- perim.north) node[fx, pos=0.6] {F1};
\draw[pub] (analyst) -- (analyst |- perim.north);
\draw[pub] (perim.south -| app) -- (app.north);
\draw[pub] (perim.south -| api) -- (api.north);
% Le tableau de bord ne joint pas l'API directement : son appel ressort par le repartiteur
% (deuxieme passage du filtrage applicatif). Le lien direct de la version anterieure
% contredisait le tableau des chemins reseau. Deux coordonnees nommees plutot qu'un
% « -| » decale : la forme (A |- B) n'accepte pas de modificateur de position.
\coordinate (dashdn) at ([xshift=-5mm]dash.north);
\coordinate (dashup) at ([xshift=3mm]dash.north);
\draw[pub] (dashdn |- perim.south) -- (dashdn);
\draw[pub] (dashup) -- (dashup |- perim.south)
  node[etq, midway, right=1pt, align=left] {JWT ;\\ 2\textsuperscript{e} passage};
\draw[priv] (app.south) -- (app.south |- net.north) node[fx, pos=0.85] {F2};
\draw[priv] (api.south) -- (api.south |- net.north);
\draw[priv] (net.south -| db) -- (db.north) node[fx, pos=0.5] {F2};
% Canal horizontal remonte de 5,05 a 5,20 : a 5,05 le trait longeait le bord
% superieur de la boite « Reseau prive » et son etiquette recouvrait le texte.
% L'etiquette est deportee sur le segment vertical (x=9,6), dans la zone libre
% a droite de la boite, ou rien ne la concurrence.
\draw[man] ([xshift=4mm]api.south) |- (9.6,5.2) -- (9.6,3.55) -- (5.4,3.55)
  -- (dw.north);
\node[etq, anchor=west] at (9.75,4.35) {identité seule (É11)};
\draw[pub] (repo) -- (ci) node[etq, midway, right] {OIDC};
\draw[pub] (ci) -- (reg) node[fx, pos=0.5] {F3};
% Canal remonte a 5,50 : les deux canaux horizontaux de cette bande (E11 a 5,20
% et celui-ci) sont ainsi separes de 0,30 et ne se croisent plus.
\draw[man] (reg.south) -- (13.0,5.5) -- (8.8,5.5) -- (dash.south);
\node[etq] at (10.7,5.5) {étiquette SHA (F3)};
\draw[jour] (perim.west) -- (-0.9,8.0);
\draw[jour] (app.west)   -- (-0.9,6.2);
\draw[jour] (net.west)   -- (-0.9,4.4);
\draw[jour] (db.west)    -- (-0.9,2.6);
\draw[jour] (-0.9,8.0) -- (-0.9,0.8) -- (col.west);
\draw[jour] (col.north -| dw) -- (dw.south) node[fx, pos=0.5] {F4};
\draw[priv, {Latex[length=1.8mm]}-{Latex[length=1.8mm]}] (dw) -- (sched) node[fx, pos=0.5] {F5};
% L'ecart entre les deux boites (1 cm) est plus etroit que le mot « connecteur » :
% l'etiquette est donc descendue sous la rangee plutot que placee dans l'entre-deux.
\draw[priv] (sched) -- (enc) node[fx, pos=0.5] {F5};
% --- separation des plans : bandes de fond, calque « backgrounds », sans surcout de hauteur
\begin{scope}[on background layer]
  \node[bandecontrole,    fit=(perim)]                 {};
  \node[bandedonnee,      fit=(app)(api)(dash)(reg)]   {};
  \node[bandecontrole,    fit=(net)]                   {};
  \node[bandedonnee,      fit=(db)(dw)(sched)(enc)]    {};
  \node[bandeobservation, fit=(col)]                   {};
\end{scope}
\node[legendecadre, anchor=north west, font=\scriptsize, text width=15.0cm] at (-0.95,-0.10)
  {\legendetrait{fluxdonnee}{chemin public}\quad
   \legendetrait{{fluxdonnee, line width=1.3pt}}{chemin réseau privé}\quad
   \legendetrait{fluxevenement}{journal (F4)}\quad
   \legendetrait{fluxcontrole}{chemin managé, identité seule (É11)}\\[2pt]
   Fond des bandes : clair = plan de données, moyen = plan de contrôle,
   sombre = plan d'observation};
\node[etq, align=center] at (11.2,1.45) {par l'interface de sortie};
\end{tikzpicture}
\caption[Vue de déploiement : composants, liens réseau et flux F1 à F5]{Vue de déploiement :
composants, liens réseau et flux F1 à F5. Le tableau de bord n'atteint les données que par l'API,
et son appel ressort par le répartiteur : il subit un second passage du filtrage applicatif.}
\label{fig:deploiement}
\end{figure}

\subsection{Topologie réseau : l'état antérieur et la cible en service}
\label{subsec:reseau-etat-cible}

\noindent\textbf{L'état antérieur, dossier d'instruction de D15.}

Les deux vues précédentes nomment le réseau sans le décrire ; la \cref{fig:topologie-reseau} le
décrit plage par plage et règle par règle, tel qu'il se lisait dans le dépôt \textbf{avant la
refonte du 11 au 25/08/2026} --- chaque affirmation reste recoupable avec le code de cette période.
Cette vue n'est pas un état des lieux mais le dossier d'instruction de D15. Une question le
résumait : \emph{par où sortait le trafic de la tâche d'enrichissement ?} Il sortait
\emph{intégralement} par le connecteur, y compris vers les interfaces publiques du fournisseur ---
seul réglage permettant à l'appel d'atteindre un service à entrée interne ; le refus général en
entrée, journalisé à la priorité la plus basse, alimentait la détection de mouvement latéral.

\begin{figure}[htbp]
\centering
\begin{tikzpicture}[
  x=1cm, y=1cm, font=\scriptsize,
  proc/.style={draw=black!70, rounded corners=3pt, fill=black!4,
               align=center, inner sep=2.5pt, minimum height=0.62cm},
  mag/.style={draw=black!70, double, double distance=0.7pt, fill=white,
              align=center, inner sep=2.5pt, minimum height=0.58cm},
  ext/.style={draw=black!70, fill=black!10, align=center,
              inner sep=2.5pt, minimum height=0.58cm},
  vide/.style={draw=black!70, densely dashed, fill=white, align=center,
               inner sep=2.5pt, minimum height=0.58cm},
  cadre/.style={draw=black!45, densely dotted, rounded corners=4pt, line width=0.8pt},
  note/.style={draw=black!45, rounded corners=2pt, fill=black!2, align=left, inner sep=4pt},
  regle/.style={note, fill=menalplanB},
  a/.style={fluxdonnee},
  ae/.style={fluxdonnee, line width=1.3pt},
  ai/.style={lienstructure, densely dashed},
  et/.style={font=\scriptsize, text=black, fill=white, inner sep=1pt, align=center},
  etg/.style={font=\scriptsize, text=black, fill=white, inner sep=1pt, align=left}]

% ===================== cadre : le reseau virtuel =============================
\draw[cadre] (0,7.40) rectangle (11.40,13.90);
\node[et, text width=5.4cm, anchor=north, align=center,
      font=\scriptsize\itshape, inner sep=2pt] at (5.85,13.80)
   {Réseau virtuel en mode personnalisé\\ un seul, partagé par tous les locataires};

% ===================== hors cadre, en haut : les charges L3 ==================
\node[ext, text width=4.0cm] (l3a) at (2.40,16.00)
   {\textbf{Cinq services et trois tâches (L3)}};
\node[ext, text width=4.0cm] (l3b) at (9.40,16.00)
   {dont la \textbf{tâche planifiée\\ d'enrichissement}};

% ===================== interieur du cadre ====================================
\node[proc, text width=10.6cm] (conn) at (5.70,12.50)
   {\textbf{Connecteur d'accès sans serveur} \quad \texttt{10.0.3.0/28} --- 2 à 10 instances\\
    \textbf{seul occupant réel} : point de passage unique de sept charges de travail
    --- sa perte les coupe simultanément};

\node[vide, text width=4.2cm] (pub) at (2.50,10.20)
   {Sous-réseau public \texttt{10.0.1.0/24}\\ \textbf{vide} \quad (É2)};
\node[vide, text width=4.2cm] (priv) at (2.50,8.40)
   {Sous-réseau privé \texttt{10.0.2.0/24}\\ \textbf{vide}, accès privé actif,\\
    journaux de flux 0,5 \quad (É2)};

\node[proc, text width=3.8cm] (app) at (8.50,10.20)
   {Plage d'appairage de services privés\\ allouée automatiquement};
\node[mag, text width=3.8cm] (bdd) at (8.50,8.40)
   {Base de données\\ adresse privée exclusive};

% ===================== hors cadre, en bas ====================================
\node[proc, text width=4.2cm] (pass) at (2.50,6.00)
   {Passerelle de sortie et routeur\\ rattachée au \textbf{seul sous-réseau privé}\\
    journalisation limitée aux erreurs};
\node[mag, text width=5.0cm] (four) at (8.70,6.00)
   {Interfaces du fournisseur\\ (entrepôt, secrets, service d'encodage)};

% ===================== a droite : les regles =================================
\node[regle, text width=3.5cm] at (13.65,11.35)
   {\textbf{Règles de filtrage en entrée}\\[2pt]
    \texttt{800} sondes de disponibilité (étiquette cible jamais attribuée)\\[2pt]
    \texttt{900} autorisation interne \texttt{10.0.1.0/24} $\leftrightarrow$
    \texttt{10.0.2.0/24}, tous ports TCP et UDP, ICMP\\[2pt]
    \texttt{1000} HTTPS \texttt{0.0.0.0/0} (étiquette cible jamais attribuée)\\[2pt]
    \texttt{65534} refus général, journalisé};
\node[regle, text width=3.5cm] at (13.65,8.10)
   {\textbf{Règles de sortie}\\[2pt] \textbf{aucune} --- rien, dans le réseau,
    ne borne ce qui sort};

% ===================== arcs ==================================================
% 1 : quatre services sur cinq
\draw[a] (2.00,15.55) -- (2.00,13.30);
\node[etg, text width=2.7cm, anchor=west] at (2.15,14.85)
   {sortie limitée aux plages privées (quatre services sur cinq)};
% 2 : la tache d'enrichissement, sortie integrale
\draw[ae] (9.60,15.55) -- (9.60,13.30);
\node[etg, text width=3.4cm, anchor=east] at (9.45,14.85)
   {\textbf{sortie intégrale} --- seul réglage permettant d'atteindre un service à entrée interne};
% 3 : connecteur -> appairage -> base
\draw[a] (8.50,11.70) -- (app.north);
\draw[a] (app.south) -- node[et, pos=0.5, xshift=6.5mm] {TCP privé} (bdd.north);
% 4 : connecteur -> interfaces du fournisseur, par l'acces prive
\draw[ae] (10.90,11.70) -- (10.90,6.65);
\node[et, anchor=west] at (11.05,7.05) {par l'accès privé};
% 5 : le sous-reseau prive est le seul desservi par la passerelle
\draw[a] (2.00,7.65) -- (2.00,6.65);
% 6 : chemin non atteste --- le connecteur n'est pas desservi
\draw[ai] (5.50,11.70) -- (5.50,5.00) -- node[croixrefus]{} (2.50,5.00) -- (2.50,5.32);

% ===================== cartouches =============================================
\node[et, text width=6.4cm, anchor=north, align=left] at (4.00,4.60)
   {le connecteur n'est pas dans la liste des sous-réseaux servis :
    \textbf{chemin non attesté par le code}};

\node[legendecadre, anchor=north east, font=\scriptsize, text width=6.4cm] at (15.60,4.95)
   {\legendetrait{fluxdonnee}{flux réseau}\quad
    \legendetrait{{fluxdonnee, line width=1.3pt}}{sortie effective}\\[2pt]
    \legendecroix{arc absent : chemin que le code n'atteste pas}};

\end{tikzpicture}
\caption[Topologie réseau avant la refonte]{Topologie réseau \textbf{avant la refonte} --- état du
dépôt jusqu'au 11/08/2026. La sortie des charges de travail n'y était pas gouvernée par le réseau
mais par le paramètre de la plateforme d'exécution : la passerelle ne desservait que le sous-réseau
privé, lequel était vide, et la plage qui portait la totalité de la sortie ne lui était pas
rattachée. Aucune règle de sortie n'existait : c'est le diagnostic qui a motivé D15.}
\label{fig:topologie-reseau}
\end{figure}

\medskip
\noindent\textbf{La refonte appliquée (D15).}
\label{subsec:reseau-cible}
La topologie qui précède décrit l'état antérieur ; celle qui suit décrit la refonte qui l'a
remplacé --- décision D15, conduite en douze étapes \textbf{du 11 au 25/08/2026}. La date de référence du
mémoire lui étant postérieure, \textbf{le réseau décrit ici est celui qui est en service à
\dateref}.

\medskip
\noindent\textbf{Le verrou que la refonte a levé.}
Six des sept charges desservies par le connecteur portaient le réglage de sortie \emph{plages
privées seulement} : seul le trafic destiné aux plages privées normalisées entrait dans le réseau
virtuel. Tout le reste --- appels aux interfaces du fournisseur, appel du tableau de bord vers le
répartiteur --- empruntait le chemin managé de la plateforme d'exécution, \textbf{hors du réseau
virtuel, donc hors de portée de toute règle de pare-feu}. Poser alors un refus par défaut en sortie
n'aurait rien bloqué : \textbf{le défaut n'était pas l'absence d'une règle, c'était que le trafic
ne passait pas là où elle se serait appliquée} --- ce qui explique aussi que la passerelle de
traduction d'adresses ne manquait à personne. La migration a donc d'abord déplacé le trafic, et
seulement ensuite l'a contrôlé.

La \emph{sortie réseau directe}, disponible en version générale sur la plateforme retenue, fait les
deux : elle place les interfaces de sortie des charges de travail \emph{dans} un sous-réseau, ce
qui rend le pare-feu applicable à elles, et elle supprime le connecteur d'accès sans serveur ---
son coût plancher permanent et le point de défaillance unique qu'il constituait pour les sept
charges qu'il desservait. Le schéma du fournisseur d'infrastructure, dans la version verrouillée
par le dépôt, expose le mécanisme sur le service comme sur la tâche : relevé du 29/08/2026.

\medskip
\noindent\textbf{Principe directeur et plan d'adressage.}
Le principe est la traduction directe du principe d'identité (PR2) au plan réseau :
\textbf{l'identité, plan de contrôle du socle, devient le sélecteur du pare-feu lui-même}. Le
critère d'acceptation qui en découle est vérifiable en une commande : \emph{aucune règle
d'autorisation ne comporte de plage source}. Le plan d'adressage qui en résulte --- quatre
sous-réseaux d'un préfixe \texttt{/26} dans un même bloc \texttt{/24}, un par locataire, le
troisième créé vide et réservé --- est publié avec la justification de chaque longueur de préfixe
au \cref{tab:adressage-cible} : un préfixe posé sans raison est une dette d'adressage.


\medskip
\noindent\textbf{Les règles de pare-feu en service.}
Douze règles --- \textbf{cinq autorisations et sept refus} --- ont remplacé les quatre antérieures,
dont une ouvrait tous les ports TCP et UDP ainsi qu'ICMP entre deux plages, sans aucune cible.
Elles ont été posées du 22 au 23/08/2026, autorisations d'abord, refus nommés ensuite, refus par
défaut en dernier ; \textbf{chaque autorisation porte une identité nommée et un port nommé}. La
journalisation suit la convention du projet --- \emph{les refus sont journalisés, les autorisations
ne le sont pas} --- pour ne pas payer un volume d'ingestion sans valeur de détection (BNF3) ; chaque
refus journalisé alimente l'export existant vers l'entrepôt, sans modifier la chaîne de détection.
Le \cref{tab:regles-cibles} les donne règle par règle --- référence, sens, identité ciblée,
destination, ports et action.


Trois lectures méritent la soutenance. \textbf{Les refus nommés E3, E7, E8, I1 et I2 ne changent
aucun verdict} --- E9 les rendrait redondants : ils existent pour \emph{produire un signal}, une
tentative du tableau de bord vers la base étant alors journalisée sous un nom de règle qui dit
exactement ce qui s'est passé, directement exploitable par la chaîne de détection. \textbf{Aucune
autorisation en entrée ne subsiste} : le pare-feu est à états, rien n'initie de connexion vers une
charge de travail par le réseau, et les trois règles antérieures ont été supprimées le 24/08/2026
sans successeur --- l'inventaire des instances de calcul, vide, établit qu'elles ne protégeaient
rien. Enfin \textbf{la règle E5 borne le tableau de bord à une seule adresse au monde}, celle du
répartiteur : le réseau referme un écart que le plan d'identité laisse ouvert (É9).

La \cref{fig:topologie-cible} rassemble le plan d'adressage et ces règles en une vue unique, à
mettre en regard de la \cref{fig:topologie-reseau} : les charges de travail y figurent
\emph{à l'intérieur} des sous-réseaux, et le bloc des règles ne comporte plus une seule plage
source.

\begin{figure}[htbp]
\centering
\begin{tikzpicture}[
  x=1cm, y=1cm, font=\scriptsize,
  sn/.style={draw=black!70, rounded corners=3pt, fill=black!4, align=left,
             inner sep=4pt, text width=4.0cm},
  snvide/.style={draw=black!70, densely dashed, fill=white, align=left,
                 inner sep=4pt, text width=2.2cm},
  mag/.style={draw=black!70, double, double distance=0.7pt, fill=white,
              align=center, inner sep=3pt, text width=3.4cm},
  ext/.style={draw=black!70, fill=black!10, align=center, inner sep=3pt,
              text width=2.8cm},
  vpc/.style={draw=black!55, rounded corners=5pt, densely dotted},
  bloc/.style={draw=black!60, rounded corners=2pt, fill=white, align=left,
               inner sep=4pt, text width=4.4cm, font=\scriptsize},
  note/.style={draw=black!45, rounded corners=2pt, fill=black!2, align=left,
               inner sep=4pt, font=\scriptsize},
  a/.style={fluxdonnee},
  et/.style={font=\scriptsize, fill=white, inner sep=1pt, text=black}]

% ============ cadre du reseau virtuel =========================================
\node[vpc, minimum width=10.05cm, minimum height=7.95cm] (rv) at (5.05,6.75) {};
% Fond blanc obligatoire : sans lui, la bordure haute du cadre barrait le titre.
% C'est l'idiome deja employe par la figure precedente (style « et »).
\node[anchor=north west, font=\scriptsize\itshape, black!60,
      fill=white, inner sep=2pt] at (1.35,10.50)
     {Réseau virtuel en mode personnalisé --- mode de routage régional \emph{déclaré}};

% ============ les deux sous-reseaux de locataire ==============================
\node[sn] (mn) at (2.55,9.40)
     {\textbf{Locataire plateforme} --- \texttt{10.0.8.0/26}\\[1pt]
      API de supervision · tableau de bord · tâche d'enrichissement};
\node[sn] (el) at (7.55,9.40)
     {\textbf{Locataire hébergé} --- \texttt{10.0.8.64/26}\\[1pt]
      application hébergée · migration · contrôle d'isolation};

% refus croise entre locataires, journalise
\draw[fluxrefuse] (mn.east) -- node[croixrefus]{} (el.west);
\node[et, anchor=north, align=center] at (5.05,9.05)
     {\textbf{E7--E8}\\ \textbf{I1--I2}};

% ============ le troisieme locataire, deja adresse ============================
\node[snvide] (t3) at (1.45,7.05)
     {\textbf{Troisième}\\ \texttt{10.0.8.128/26}\\ \emph{déjà adressé}};
\node[snvide] (rs) at (8.65,7.05)
     {\textbf{Réserve}\\ \texttt{10.0.8.192/26}\\ \emph{ferme le \texttt{/24}}};

% ============ appairage et base ===============================================
\node[mag] (psa) at (5.05,4.95) {Appairage de services privés --- \texttt{/20}};
\node[mag] (bd)  at (5.05,3.65) {Base de données --- adresse privée exclusive};
\draw[a] (psa) -- (bd);

% ============ hors du reseau virtuel ==========================================
\node[ext] (api) at (13.05,10.55) {Interfaces privées du fournisseur};
\node[ext] (lb)  at (13.05,8.95)  {Répartiteur\\ \emph{une seule adresse}};

% ============ arcs vers la base, par les couloirs libres =======================
\draw[a] (3.90,8.98) -- node[et, pos=0.62] {\textbf{E1} --- 5432} (3.90,5.42);
\draw[a] (6.20,8.98) -- node[et, pos=0.62] {\textbf{E2} --- 5432} (6.20,5.42);

% ============ arcs sortants ===================================================
\draw[a] (el.east) -- node[et, pos=0.52] {\textbf{E4} --- 443} (api.west);
% routage orthogonal au-dessus du cadre : l'arc ne croise ni le titre ni un noeud
\draw[a] (0.40,9.92) -- (0.40,11.10)
         -- node[et, pos=0.5] {\textbf{E5} --- 443 (tableau de bord seul)} (11.30,11.10)
         -- (11.30,8.95) -- (lb.west);

% ============ bloc des regles =================================================
\node[bloc, anchor=north west, fill=menalplanB] (rg) at (10.75,7.55)
     {\textbf{Règles --- sortie}\\[1pt]
      \texttt{1000} base de données, par identité\\
      \texttt{1200} interfaces du fournisseur\\
      \texttt{1250} répartiteur, une seule adresse\\
      \texttt{1400} refus croisé, \emph{journalisé}\\
      \texttt{65534} \textbf{refus par défaut}, \emph{journalisé}\\[3pt]
      \textbf{Règles --- entrée}\\[1pt]
      \texttt{1400} refus croisé, \emph{journalisé}\\
      \texttt{65534} refus général, \emph{journalisé}\\
      \emph{aucune autorisation en entrée}};

% ============ ce qui disparait, et ce que la figure demontre ===================
\node[note, anchor=north west, text width=9.75cm] at (0.20,2.55)
     {\textbf{Supprimés} --- le connecteur d'accès sans serveur et sa plage
      \texttt{10.0.3.0/28}, la passerelle de traduction d'adresses et son routeur, les deux
      sous-réseaux vides \texttt{10.0.1.0/24} et \texttt{10.0.2.0/24}, et les trois règles
      d'entrée qui ciblaient une étiquette que rien ne porte.};

\node[legendecadre, anchor=north west, text width=4.4cm, font=\scriptsize] at (10.75,2.55)
     {\legendetrait{fluxdonnee}{flux autorisé, par identité}\\[2pt]
      \legendecroix{refus opposé par le pare-feu}};

\end{tikzpicture}
\caption[Topologie réseau après la refonte appliquée]{Topologie réseau \textbf{après la refonte
appliquée du 11 au 25/08/2026} --- le réseau en service à \dateref, à comparer à la
\cref{fig:topologie-reseau}. Les charges de travail émettent \emph{depuis} un sous-réseau : le
pare-feu leur est applicable, et aucune autorisation ne comporte de plage source.}
\label{fig:topologie-cible}
\end{figure}

\begin{alertbox}[title={Une limite du mécanisme, dite avant qu'un jury la trouve}]
Le sélecteur par compte de service est une fonctionnalité réelle du pare-feu du fournisseur, mais
\textbf{son application aux interfaces de sortie directe n'est pas établie} : à \dateref, le schéma
du fournisseur d'infrastructure n'expose que les étiquettes à ce niveau. La refonte applique donc
une convention --- \emph{une étiquette par compte de service, l'une et l'autre déclarées dans le
même bloc} --- et un commutateur qui bascule vers le sélecteur par identité une fois la
vérification passée : la propriété revendiquée ne change pas, c'est le mécanisme qui l'exprime qui
attend sa preuve. La conséquence porte sur la ségrégation entre locataires et vaut pour le réseau
en service : une étiquette se pose dans la définition de la révision, de sorte que \textbf{la
frontière entre locataires est opposable à une charge de travail compromise, mais non à un
opérateur capable de déployer} --- lequel poserait l'étiquette de l'autre locataire. Ce second cas
relève du resserrement du droit de déploiement (\cref{subsec:delegation}). La vérification qui lève
la réserve est la première du protocole et tient en une règle de test et un appel.
\end{alertbox}

\medskip
\noindent\textbf{Ce que la refonte a fermé, ce qu'elle coûte, et dans quel ordre.}
Trois défauts sont tombés ensemble parce que \textbf{la passerelle de traduction d'adresses n'a pas
été mieux configurée : elle est devenue sans objet} au sens de PR1
(\cref{sec:trente-trois-briques}). Sa suppression, le 24/08/2026, a emporté son rattachement au
mauvais sous-réseau, sa journalisation limitée aux erreurs --- qui laissait toute connexion
sortante réussie sans signal --- et son allocation d'adresses automatique, qui interdisait toute
liste d'autorisation chez un tiers.

\textbf{Les montants qui suivent sont des ordres de grandeur de tarification publique, non des
montants relevés sur facture.} La suppression du connecteur, facturé au plancher de deux instances
permanentes, libère de l'ordre de \textbf{onze à treize euros par mois} ; sous-réseaux, routes et
règles ne sont pas facturés ; seuls les journaux de flux deviennent un poste de un à deux euros. S'y
ajoute un coût évité qui ne figure dans aucun tableau, la montée en charge du connecteur, qui
disparaît \emph{en tant que risque}. La conséquence est structurante : \textbf{le coût marginal
réseau d'un locataire supplémentaire passe d'une douzaine d'euros par mois à zéro} --- un
sous-réseau et deux règles de refus croisé ne sont pas facturés, et le troisième sous-réseau est
déjà créé --- de sorte que \textbf{la ségrégation réseau entre locataires a cessé d'être un
arbitrage entre coût et sécurité} et que le troisième seuil du \cref{ch:validation} devient
franchissable sans renégocier le budget. \emph{Le rapprochement de ces ordres de grandeur avec la
facture appartient au mois suivant la migration, sur les postes d'accès réseau, de traduction
d'adresses et de télémétrie.}

Les douze étapes ont obéi à une règle unique --- \emph{on n'interdit jamais avant d'avoir observé,
et on n'observe qu'après avoir fait passer le trafic par l'endroit où on l'observe} --- qui commande
le calendrier du 11 au 25/08/2026 ; la conduite de l'opération, la fenêtre d'observation de sept
jours et le traitement de la seule étape réellement dangereuse, le refus par défaut en sortie, sont
rapportés en \cref{subsec:migration-reseau}. Deux propriétés sont ici établies par contrôle
direct --- l'exposition de la sortie réseau directe dans la version verrouillée du fournisseur,
relevée le 29/08/2026, et le fait que les trois règles d'entrée supprimées ne protégeaient rien,
l'inventaire des instances de calcul étant vide ; les deux autres, applicabilité du sélecteur par
identité et démonstration que le tableau de bord ne joint plus qu'une adresse, portent leur commande
et leur critère au plan de validation continue (\cref{sec:plan-validation}). Enfin, une
architecture de pare-feu à douze règles exploitée par une seule personne se dégrade au premier
incident où une autorisation large débloque une urgence : le critère d'acceptation ---
\emph{aucune règle d'autorisation ne comporte de plage source} --- est donc porté en porte
bloquante par l'analyse de configuration de la chaîne de provisionnement (\cref{ch:realisation}).
\textbf{L'architecture refuse elle-même sa propre dégradation.}

\medskip
\noindent\textbf{Ce que la refonte ne ferme pas.}
Cinq limites sont assumées. La plage d'appairage de services privés est immuable sur l'instance
existante : la cible \texttt{/20} ne vaut que pour un provisionnement neuf. La base de données et
l'entrepôt restent partagés --- \textbf{la ségrégation réseau n'empêche pas deux identités
d'atteindre le même serveur}, elle empêche l'une d'atteindre l'\emph{interface} de l'autre. Le
refus par défaut en sortie supprime le canal de commande vers un hôte arbitraire mais \textbf{non
l'exfiltration par les interfaces du fournisseur lui-même} : l'autorisation qui rend l'entrepôt
joignable rend joignable tout entrepôt du même fournisseur, et fermer ce chemin exigerait un
périmètre de service, donc une organisation cloud dont le projet ne dispose pas. Le détachement des
deux charges qui n'ont aucune dépendance privée leur laisse une \textbf{sortie managée non
gouvernée}, arbitrage dont le prix est mesuré par T14. Enfin la conception ne tient pas au-delà de
trois ou quatre locataires, le nombre de refus croisés croissant comme le carré de leur nombre :
au-delà, la réponse est \textbf{un projet cloud par locataire}, qui règle du même coup la base et
l'entrepôt partagés.

\subsection{Isolation des données entre locataires (D16)}
\label{subsec:isolation-donnees}

La ségrégation réseau appliquée en \cref{subsec:reseau-cible} empêche une charge de travail
d'atteindre l'interface d'une autre ; elle ne peut rien sur le fait que les deux atteignent
\emph{le même serveur} de base de données. La frontière de donnée reste donc portée par les droits
du moteur, appliqués à l'exécution : cette sous-section établit contre qui elle tient, puis expose
la conception qui la remplacerait --- décision D16, arrêtée le 29/08/2026 et \textbf{non
appliquée}. Ce statut est un choix : \textbf{deux migrations d'infrastructure dans la même
quinzaine, conduites par un opérateur unique, ne seraient pas soutenables} --- la contrainte
humaine du \cref{ch:cadre-existant} appliquée à l'ordonnancement du projet lui-même.

\medskip
\noindent\textbf{Ce que la frontière actuelle garantit, et contre qui elle cède.}
Dix ordres, dont six révocations, retirent aux comptes applicatifs l'appartenance au rôle
privilégié de l'instance et referment les connexions croisées. Le \cref{tab:adversaires-donnees}
confronte cette frontière à chaque adversaire plutôt qu'à une échelle de maturité : c'est la seule
façon de dire ce qu'elle vaut.

{\setlength{\tabcolsep}{4pt}
\begin{longtable}{|P{5.4cm}|P{2.5cm}|P{7.1cm}|}
\caption[Ce que la frontière d'isolation des données arrête, et ce qu'elle n'arrête pas]{Ce que
la frontière d'isolation des données arrête, et ce qu'elle n'arrête pas. Les trois défaillances
sont des \emph{circonstances d'exploitation}, non des attaques --- c'est précisément pourquoi une
frontière portée par un geste ne suffit pas.}
\label{tab:adversaires-donnees} \\
\hline
\textbf{Adversaire} & \textbf{La frontière tient-elle ?} & \textbf{Ce qui la tient réellement} \\
\hline
\endfirsthead
\caption[]{Ce que la frontière d'isolation des données arrête (suite)} \\
\hline
\textbf{Adversaire} & \textbf{La frontière tient-elle ?} & \textbf{Ce qui la tient réellement} \\
\hline
\endhead
Charge applicative compromise, dans l'un ou l'autre sens & \textbf{Oui} --- établi le 19/08/2026
par six vérifications sur six, par tentative de connexion réelle et non par lecture de
configuration & Le retrait de l'appartenance privilégiée et la révocation de connexion. Rien d'autre
\\ \hline
Fuite du mot de passe applicatif & \textbf{Oui} & La même révocation : le mot de passe n'ouvre que
sa propre base \\ \hline
Identité de la chaîne de livraison compromise & \textbf{Partiellement} : elle déploie sous
l'identité d'un locataire et lit ses secrets, jamais ceux de l'autre & \textbf{La liaison secret par
secret du coffre --- la frontière la plus solide du socle est dans le plan d'identité, non dans le
moteur de base} \\ \hline
Opérateur détenteur d'un rôle d'administration de l'instance & \textbf{Non, et aucune conception de
droits ne peut y répondre} : l'export lit n'importe quelle base sans traverser le moteur
d'autorisation & Rien. \emph{Un octroi n'est pas une défense contre le plan de contrôle qui
l'héberge} \\ \hline
Réconciliation de l'infrastructure après un clonage & \textbf{Non, et sans borne temporelle} : le
clone porte un nom neuf, la déclaration des comptes y est ancrée, l'application suivante les recrée
privilégiés & Rien \\ \hline
Accueil d'un nouveau locataire & \textbf{Non, et sans aucune détection} : le contrôle quotidien ne
connaît que deux comptes écrits en dur & Rien \\ \hline
Restauration d'un locataire & \textbf{Sans objet --- et c'est le problème} : l'unité de reprise
est l'instance, restaurer l'un fait reculer l'autre & Rien. Défaut de rayon d'action, non de
droits \\
\hline
\end{longtable}}

Deux enseignements commandent la conception. La frontière \textbf{cède dans trois circonstances
d'exploitation nommées} --- l'accueil d'un locataire, la réconciliation d'une infrastructure
décrite avec une infrastructure restaurée, et sa propre disparition silencieuse : ce ne sont pas
des attaques. Et elle est \textbf{asymétrique} : franchie depuis le locataire hébergé, elle livre
des empreintes de mots de passe et un secret de second facteur chiffré ; franchie depuis la
plateforme, elle livre \textbf{un identifiant national en clair}. La direction qui coûte part donc
de l'hébergeur, et l'actif qu'elle protège ne l'est par rien d'autre que la révocation.

\medskip
\noindent\textbf{La conception retenue : une isolation proportionnée à la sensibilité.}
Quatre niveaux ont été évalués --- durcissement logique remis sous code, séparation par schéma,
instance dédiée par locataire, projet cloud dédié par locataire. D16 retient le troisième
\emph{pour le seul locataire qui porte une donnée irréparable} et le premier partout ailleurs,
pour deux raisons. \textbf{La frontière doit d'abord changer de nature, pas de qualité} : remettre
les dix ordres sous contrôle du code les rend versionnés, rejoués et surveillés, mais ils restent
un geste --- à la question « que se passe-t-il si quelqu'un accorde le rôle privilégié ? », la
réponse demeure \emph{on le voit sous vingt-quatre heures}, quand avec une instance dédiée la
question n'a plus de sens. \textbf{La sensibilité commande ensuite la proportion, non l'égalité de
traitement} : le \cref{ch:cadre-existant} fonde le niveau d'exigence du socle sur la biométrie
vocale et l'identifiant national, et traiter les deux locataires à l'identique reviendrait à
surpayer pour la plateforme ou à sous-protéger l'application hébergée.

Le \textbf{test d'admission d'un locataire} énonce cette proportion en critère opposable, en trois
questions : \textbf{(i)} la donnée traitée est-elle \textbf{irréversible} en cas de fuite ---
identifiant officiel, biométrie, donnée de santé, moyen de paiement ? \textbf{(ii)} l'éditeur
est-il un \textbf{tiers} de l'exploitant du socle ? \textbf{(iii)} un engagement
\textbf{contractuel} d'hébergement est-il pris ? \textbf{Une seule réponse affirmative impose une
instance dédiée} ; aucune, une base dédiée sur l'instance mutualisée, derrière la tâche de
durcissement idempotente. L'application hébergée répond oui à la première ; un troisième locataire
répondrait très probablement oui à la deuxième.

\medskip
\noindent\textbf{Le coût, et l'avertissement qui va avec.}
Une instance dédiée zonale, sauvegardes et restauration à un instant donné identiques, coûte de
l'ordre de dix à douze euros par mois, quand la refonte réseau en a libéré onze à treize --- deux
ordres de grandeur de tarification publique, le second attendant son rapprochement avec la facture.
\textbf{L'économie du réseau finance exactement l'isolation des données, et une seule fois} : le
solde commun des deux décisions est proche de zéro, et c'est cette neutralité qui rend D16 tenable
sous la contrainte économique du \cref{ch:cadre-existant}. Le prix payé doit être dit : le
locataire hébergé perdrait la bascule régionale automatique, acquise le 08/08/2026 et
\textbf{jamais éprouvée par un test} --- la seule reprise mesurée du projet est un clonage, qui
fonctionne identiquement sur une instance zonale --- et \textbf{la mesure de 32~min~45~s ne se
transporte pas}, ayant été relevée sur une instance donnée. En contrepartie, la nouvelle instance
naîtrait avec le chiffrement par clé gérée, à coût nul : possibilité que l'instance existante a
perdue définitivement, le champ n'étant modifiable qu'à la création.

\medskip
\noindent\textbf{Le piège de la reprise, et pourquoi le contrôle ne doit pas réparer.}
Puisque la frontière ne survit pas à la réconciliation qui suit un clonage, elle doit être déplacée
dans la définition même de \emph{reprise terminée} : le mode opératoire de restauration ne s'achève
pas au retour du service, mais au rejeu de la tâche de durcissement puis du contrôle d'isolation
--- motif déjà éprouvé par la tâche de migration du locataire. S'y ajoutent une alerte de
fraîcheur, l'absence du contrôle étant aujourd'hui silencieuse, et l'énumération des comptes plutôt
que leur inscription en dur. Une réparation automatique est en revanche \textbf{écartée par un
argument dirimant} : réparer exige précisément les privilèges que la frontière retire, et un
réparateur permanent supposerait de \textbf{créer une identité inter-bases à demeure pour défendre
une frontière inter-bases}. La remise sous code passe pour la même raison par une tâche idempotente
et non par un fournisseur de description supplémentaire, l'instance n'ayant aucune adresse publique
et la chaîne de provisionnement ne contenant ni planification ni application. Son gain réel n'est
d'ailleurs pas l'écriture des ordres, c'est que \textbf{la source de vérité qui crée la base
devienne celle qui en referme la porte}.

\medskip
\noindent\textbf{Ce que la conception ne couvre pas.}
Elle ne prétend rien contre le plan de contrôle du fournisseur : un opérateur détenteur d'un rôle
d'administration exporte n'importe quelle base sans traverser le moteur d'autorisation, et seule
une séparation de projets y répondrait --- réponse terminale déjà nommée en
\cref{subsec:reseau-cible}. Elle ne referme pas le chiffrement applicatif de l'identifiant
national, chantier propre à l'éditeur. Enfin elle suppose établie une propriété qui ne l'est pas :
la version de moteur en service autorise un compte portant le droit de création de rôles à
s'accorder lui-même n'importe quelle appartenance, comportement durci seulement à la version
suivante. \textbf{Tant que l'énumération des attributs de rôle n'a pas été relevée, la garantie
établie le 19/08/2026 porte une réserve} --- et c'est la première vérification du protocole, celle
qui conditionne les autres.

\medskip
\noindent\textbf{Critères de validation de l'état final.}
Une vue d'architecture ne vaut que si un tiers peut dire à quoi elle sert et comment il en prononce
la recette. Le \cref{tab:recette-vues} le fait pour les \textbf{six vues qui portent une décision
opposable} : les quatre de cette section, le plan d'identité (\cref{sec:plan-identite}) et les
frontières de confiance du \cref{ch:besoins-menaces}. Les autres figures du mémoire n'y sont pas inscrites, et le
motif est toujours le même : on ne prononce pas la recette d'un état qu'on a supprimé
(\cref{fig:topologie-reseau}) ; une vue de latence ou de schéma se valide par une mesure ou par une
propriété du schéma, non par une revue d'architecture (\cref{fig:cycle-donnee},
\cref{fig:modele-donnees}) ; les vues de comportement illustrent un protocole sans le remplacer
(\cref{fig:seq-refus}, \cref{fig:seq-mfa}) ; et la chaîne de délégation
(\cref{fig:delegation-identite}) \textbf{n'est pas recettable par construction}, puisqu'elle décrit
ce qui serait atteignable et non ce qui est vérifié --- sa réserve est son objet entier.

\begin{table}[htbp]
\caption[Critères de recette des vues d'architecture]{Critères de recette des vues d'architecture :
le point de vue servi, les critères datés avec leur protocole, et la réserve de chaque vue.}
\label{tab:recette-vues}
\centering
\footnotesize
\setlength{\tabcolsep}{4pt}
\renewcommand{\arraystretch}{1.05}
\begin{tabular}{|P{2.15cm}|P{3.20cm}|P{5.35cm}|P{4.00cm}|}
\hline
\textbf{Vue} & \textbf{Point de vue --- décision servie} & \textbf{Critères de recette, datés, et
protocole} & \textbf{Réserve --- ce que la vue ne démontre pas} \\
\hline
Contexte & Architecte --- arbitre les dépendances non maîtrisées & Les trois imports différés ne
sont sur aucun chemin de requête (conception, \dateref) ; la zone de noms reste hors code ---
écart É10 & Ni volumétrie, ni disponibilité des tiers \\ \hline
Couches & Architecte --- vérifie qu'aucune couche n'est traversée sans identité & L2 vérifiée à
chaque saut : \textbf{T5} conforme 23/08/2026 ; \textbf{T6} partiel au \dateref & Ne montre pas les
chemins réseau ; voir la vue de déploiement \\ \hline
Déploiement & Exploitant --- décide si le chemin privé est le seul chemin de données & Sorties
refusées : \textbf{T20} conforme 24/08/2026 ; données inaccessibles depuis Internet : \textbf{T4}
partiel & Ni volumes, ni latences, ni perte du connecteur (\textbf{T15}, validation continue) \\
\hline
Topologie cible & Architecte réseau --- prononce la recette de la refonte & Émission depuis un
sous-réseau, pare-feu applicable : \textbf{T20} conforme 24/08/2026 ; isolation : \textbf{T5}
conforme 23/08/2026 & Ne couvre pas le trafic entrant ; \textbf{T1} reste non conforme \\ \hline
Plan d'identité & Responsable sécurité --- décide si une identité peut altérer une preuve &
Inventaire de clés vide : \textbf{T7} conforme 19/08/2026 ; intégrité de la preuve : \textbf{T11}
conforme, test automatisé & Décrit les droits déclarés, non leur composition (voir la vue de
délégation) \\ \hline
Frontières de confiance & Analyste de la menace --- décide où poser un contrôle & TB1--TB4
franchies par le réseau, TB5 par l'identité seule : \textbf{T5} et \textbf{T11} conformes ;
\textbf{T14} non conforme au \dateref & Ne hiérarchise pas les scénarios ; la gravité est portée
par l'annexe~\ref{ann:modele-menaces} \\
\hline
\end{tabular}
\end{table}

\section{Vue dynamique : les sept flux et le cycle de la donnée}

La vue en couches décrit ce qui existe ; la vue en flux décrit ce qui se passe. Les sept flux F1
à F7 forment le vocabulaire commun du projet, employé pour l'analyse de menace du \cref{ch:besoins-menaces}, la
réalisation et les tests (\cref{tab:flux}).

\begin{table}[htbp]
\caption{Les sept flux de l'architecture.}
\label{tab:flux}
\centering
\footnotesize
\setlength{\tabcolsep}{3.5pt}
\renewcommand{\arraystretch}{1.05}
\begin{tabular}{|P{1.1cm}|P{10.2cm}|P{3.8cm}|}
\hline
Flux & Chemin & Exigences \\
\hline
F1 & Utilisateur $\to$ périmètre $\to$ filtrage $\to$ charge de travail & EX1--EX3 \\ \hline
F2 & Charge de travail $\to$ réseau privé $\to$ base de données & EX4--EX6 \\ \hline
F3 & Dépôt $\to$ chaîne de livraison $\to$ registre $\to$ déploiement & EX7--EX10 \\ \hline
F4 & Toutes les couches $\to$ collecte $\to$ entrepôt de supervision & EX11--EX13 \\ \hline
F5 & Planificateur $\to$ enrichissement $\to$ modèle $\to$ recherche vectorielle & EX11, EX14, EX15 \\ \hline
F6 & Vulnérabilités de livraison $\to$ encodage $\to$ rapprochement $\to$ priorisation & apport, aucune exigence \\ \hline
F7 & Dépôt $\to$ plan $\to$ contrôle $\to$ approbation $\to$ application & EX16--EX18 \\
\hline
\end{tabular}
\end{table}


\textbf{Le flux F6 est l'hypothèse d'apport de l'architecture.} Il relie deux ensembles
habituellement disjoints : la chaîne de livraison, qui classe les vulnérabilités par gravité
théorique, et la chaîne de détection, qui observe les techniques réellement dirigées contre le
système. Le rapprochement réordonne les correctifs --- une vulnérabilité moyenne liée à une
technique observée cette semaine passe devant une vulnérabilité élevée sans rapport avec l'activité
constatée. F6 ne couvre aucune exigence et ne réduit aucun risque par lui-même : c'est une
\emph{hypothèse}, publiée avec le protocole qui permettrait de la réfuter --- la comparaison des
ordres de correction obtenus avec et sans rapprochement. Ce protocole ne figure pas parmi les six du
plan de validation continue ; il est écrit pour qu'un tiers l'exécute et contredise l'hypothèse.

\subsection{Cycle de vie de la donnée de sécurité}

Les flux disent qui parle à qui, non ce que devient la donnée. Six étapes séparent l'événement de
son affichage, et ce sont les \emph{cadences portées sur les transitions} qu'il faut lire dans la
\cref{fig:cycle-donnee} : cinq minutes pour la normalisation, cinq pour la détection, quinze pour
l'enrichissement sont des attentes d'ordonnanceur et non des temps de traitement, qui s'additionnent
avec la position de l'événement dans la fenêtre et bornent le délai bien plus sûrement que la
puissance de calcul. Deux propriétés s'y lisent aussi : le score est recalculé à chaque consultation
et jamais persisté ; et la branche d'enrichissement est \emph{informationnelle}, la
\cref{fig:modele-donnees} en donnant la raison, de modèle de données et non d'algorithme.

\begin{figure}[htbp]
\centering
\begin{tikzpicture}[
  x=1cm, y=1cm, font=\scriptsize,
  proc/.style={draw=black!70, rounded corners=3pt, fill=black!4,
               align=center, inner sep=3pt, minimum height=0.62cm},
  lat/.style={draw=black!70, rounded corners=3pt, fill=white,
              align=center, inner sep=3pt, minimum height=0.62cm},
  num/.style={circle, draw=black!70, fill=white, inner sep=0.8pt,
              minimum size=0.42cm, font=\scriptsize\bfseries},
  note/.style={draw=black!45, rounded corners=2pt, fill=black!2, align=left, inner sep=4pt},
  a/.style={fluxdonnee},
  ad/.style={fluxevenement},
  nul/.style={lienstructure, densely dashed},
  et/.style={font=\scriptsize, text=black, fill=white, inner sep=1.5pt, align=center},
  etg/.style={font=\scriptsize, text=black, fill=white, inner sep=1.5pt, align=left}]

% ===================== chaine principale, verticale ==========================
\node[proc, text width=5.8cm] (b1) at (4.00,14.60)
   {\textbf{Ingestion} --- quatre exports : charges de travail,\\
    verdicts du filtrage applicatif, refus du pare-feu,\\ base de données};
\node[proc, text width=5.8cm] (b2) at (4.00,12.20)
   {\textbf{Normalisation} --- cinq requêtes planifiées,\\
    fenêtre glissante 15~min, fusion anti-doublon};
\node[proc, text width=5.8cm] (b3) at (4.00,9.70)
   {\textbf{Détection} --- sept requêtes planifiées R1--R7,\\
    fenêtre 15~min, insertion conditionnelle};
\node[proc, text width=5.8cm] (b5) at (4.00,6.20)
   {\textbf{Scoring} --- recalculé à la lecture, \textbf{jamais persisté} ;\\
    poids par gravité et prime de chaîne d'attaque\\[1pt]
    \emph{le score n'existe qu'en mémoire :\\ l'ordre d'hier n'est pas reconstituable}};
\node[proc, text width=5.8cm] (b6) at (4.00,3.30)
   {\textbf{Affichage} --- API de supervision,\\ rendu serveur, navigateur};

% ===================== derivation laterale ===================================
\node[lat, text width=5.0cm] (b4) at (11.80,9.70)
   {\textbf{Enrichissement} --- encodage puis recherche vectorielle,\\
    trois candidats, seuil 0,60\\[1pt]
    \emph{branche informationnelle}};

% ===================== pastilles de numerotation =============================
\node[num, anchor=east] at ([xshift=-2.5pt]b1.west) {1};
\node[num, anchor=east] at ([xshift=-2.5pt]b2.west) {2};
\node[num, anchor=east] at ([xshift=-2.5pt]b3.west) {3};
\node[num, anchor=south east] at ([shift={(-2.5pt,1.5pt)}]b4.north west) {4};
\node[num, anchor=east] at ([xshift=-2.5pt]b5.west) {5};
\node[num, anchor=east] at ([xshift=-2.5pt]b6.west) {6};

% ===================== horodatage relatif ====================================
% Repris de l'ancienne chronologie de l'incident : l'origine t0 est le refus
% oppose au perimetre, avant toute charge de travail.
\node[etg, anchor=west] at (7.10,14.60) {$t_0$ + qq.~s};
\node[etg, anchor=west] at (7.10,12.20) {$t_0$ + 0 à 5~min};
\node[etg, anchor=west] at (7.10,10.35) {$t_0 + \approx 15$~min};
\node[etg, anchor=west] at (7.10,6.20)  {à la consultation};
\node[etg, anchor=west] at (7.10,3.30)  {+ 0 à 10~s};

% ===================== arcs : les cadences sont sur les fleches ==============
\draw[a] (b1.south) -- node[et, pos=0.5] {\textbf{quasi temps réel}} (b2.north);
\draw[a] (b2.south) -- node[et, pos=0.5] {\textbf{5~min} (cadence)} (b3.north);
\draw[ad] (b3.east) -- node[et, pos=0.5] {\textbf{15~min}} (b4.west);
\draw[a] (b3.south) -- node[et, pos=0.35] {\textbf{à la requête}} (b5.north);
\draw[a] (b5.south) -- node[et, pos=0.5] {\textbf{10~s} (rafraîchissement)} (b6.north);

% ===================== l'arc qui n'existe pas ================================
\draw[nul] (b4.south) -- node[croixrefus]{} (b5.east);
\node[etg, anchor=north west] at (10.10,7.10)
   {\textbf{arc absent} --- pas de clé d'entité :\\ jointure impossible\\
    (voir le modèle de données)};

% ===================== legende interne ========================================
\node[legendecadre, anchor=north west, text width=14.9cm, font=\scriptsize] at (0.60,2.35)
   {\legendetrait{fluxdonnee}{passage d'étape}\quad
    \legendetrait{fluxevenement}{branche d'enrichissement}\quad
    \legendecroix{arc absent}\\[2pt]
    Origine $t_0$ : le refus opposé au périmètre (L1), avant toute charge de travail.
    Les horodatages relatifs décrivent la \emph{configuration} ; seuls les repères de
    détection ont été mesurés (\cref{sec:bout-en-bout}).};

\end{tikzpicture}
\caption[Cycle de vie de la donnée de sécurité]{Cycle de vie de la donnée de sécurité : six
étapes, le budget de latence porté par les cadences et l'horodatage relatif de chaque étape. Le
délai entre l'événement et son affichage est fait d'\emph{attentes d'ordonnanceur} et non de temps
de traitement : trois attentes s'additionnent (5 + 5 + 15~min), et aucune puissance de calcul ne
les réduirait. La branche d'enrichissement n'alimente pas le calcul de score --- impasse
fonctionnelle assumée.}
\label{fig:cycle-donnee}
\end{figure}

\subsection{Modèle de données des deux systèmes de stockage}

Le socle stocke dans deux systèmes de nature opposée : une base relationnelle pour l'exploitation
du tableau de bord, un entrepôt analytique pour les preuves de sécurité. Dans le modèle
entité-association de la \cref{fig:modele-donnees}, deux absences comptent plus que les présences.
La table d'enrichissement porte l'identifiant de la détection dont elle est issue, mais pas
l'entité que cette détection désigne, or le score d'incident regroupe les détections \emph{par
entité} : la jointure du rattachement sémantique au score est structurellement impossible,
conséquence du schéma et non choix d'algorithme. Et aucun lien n'existe entre les deux systèmes,
la chaîne de preuve étant cloisonnée de la base des comptes --- propriété voulue dont le prix est
qu'un verdict ne se résout pas en un nom sans franchir volontairement ce cloisonnement. Le détail
des dix tables, avec leur écrivain unique, figure au \cref{ch:realisation}
(\cref{tab:modele-donnees}).

\begin{figure}[htbp]
\centering
\begin{tikzpicture}[
  x=1cm, y=1cm, font=\scriptsize,
  eh/.style={draw=black!70, fill=white, align=left, inner sep=3pt, anchor=north,
             font=\scriptsize\ttfamily\bfseries},
  eb/.style={draw=black!70, fill=white, align=left, inner sep=3pt, anchor=north},
  ehg/.style={draw=black!70, fill=black!10, align=left, inner sep=3pt, anchor=north,
              font=\scriptsize\ttfamily\bfseries},
  ebg/.style={draw=black!70, fill=black!10, align=left, inner sep=3pt, anchor=north},
  trait/.style={draw=black!70, rounded corners=3pt, fill=black!4, align=center,
                inner sep=3pt, anchor=north},
  cadre/.style={draw=black!45, densely dotted, rounded corners=4pt, line width=0.8pt},
  note/.style={draw=black!45, rounded corners=2pt, fill=black!2, align=left, inner sep=4pt},
  rel/.style={black!70, line width=0.7pt},
  relg/.style={black!55, line width=0.7pt},
  a/.style={fluxdonnee},
  et/.style={font=\scriptsize, text=black, fill=white, inner sep=1.5pt, align=center},
  etg/.style={font=\scriptsize, text=black, fill=white, inner sep=1.5pt, align=left},
  pics/patte/.style={code={\draw (0,0) -- (0.30,0.17); \draw (0,0) -- (0.30,0);
                           \draw (0,0) -- (0.30,-0.17);}},
  pics/un/.style={code={\draw (0,-0.15) -- (0,0.15);}}]

% ============================ legende de notation ============================
\node[anchor=north west, align=left, text width=15.4cm, inner sep=0pt,
      font=\scriptsize] at (0.00,18.40)
  {\textbf{Notation patte-d'oie} \quad
   \tikz[baseline=-0.5ex]{\draw[rel](0,0)--(0.55,0);\draw[rel](0.40,-0.15)--(0.40,0.15);}\;un
   \quad
   \tikz[baseline=-0.5ex]{\draw[rel](0,0)--(0.35,0);\draw[rel](0.35,0)--(0.65,0.17);
     \draw[rel](0.35,0)--(0.65,0);\draw[rel](0.35,0)--(0.65,-0.17);}\;plusieurs
   \quad
   \tikz[baseline=-0.5ex]{\draw[rel,densely dashed](0,0)--(0.9,0);}\;rattachement sans clé
   \quad
   \tikz[baseline=-0.5ex]{\draw[relg,densely dashed](0,0)--(1.0,0)
     node[croixrefus, pos=0.5]{};}\;jointure impossible\\
   \tikz[baseline=-0.5ex]{\node[ehg,anchor=center,minimum width=0.55cm,
     minimum height=0.24cm,inner sep=1pt]{};}\;table grisée : \emph{déclarée, jamais
   utilisée} --- le grisé est toujours doublé par l'étiquette écrite dans la table};

% ================================ CADRE A ====================================
\draw[cadre] (0,12.95) rectangle (15.40,17.20);
\node[anchor=north west, font=\scriptsize\bfseries] at (0.15,17.14)
  {Cadre A --- base relationnelle de l'application de supervision};

\node[eh, text width=2.4cm] (roh) at (1.55,16.15) {roles};
\node[eb, text width=2.4cm] (rob) at (roh.south)
  {\texttt{id} \textbf{PK} \,\textperiodcentered\, \texttt{name} \textbf{UNIQUE}};

\node[eh, text width=4.2cm] (ush) at (5.90,16.50) {users};
\node[eb, text width=4.2cm] (usb) at (ush.south)
  {\texttt{id} \textbf{PK} \,\textperiodcentered\, \texttt{email} \textbf{UNIQUE}
   \,\textperiodcentered\, \texttt{role\_id} \textbf{FK}\\
   mot de passe haché \,\textperiodcentered\, secret du second facteur \emph{chiffré}};

\node[eh, text width=3.8cm] (auh) at (11.60,16.50) {audit\_logs};
\node[eb, text width=3.8cm] (aub) at (auh.south)
  {\texttt{id} \textbf{PK} \,\textperiodcentered\, \texttt{user\_id} \textbf{FK} \emph{nullable}\\
   action \,\textperiodcentered\, ressource \,\textperiodcentered\, adresse
   \,\textperiodcentered\, statut \,\textperiodcentered\, horodatage \emph{indexé}};

\node[ehg, text width=6.6cm] (aph) at (5.90,14.15) {api\_keys};
\node[ebg, text width=6.6cm] (apb) at (aph.south)
  {empreinte de clé \textbf{UNIQUE} \,\textperiodcentered\, \texttt{user\_id} \textbf{FK}
   \,\textperiodcentered\, \textbf{déclarée, inerte} : aucun mécanisme d'authentification par clé};

% ---- relations du cadre A ----
\draw[rel] (2.806,15.45) -- (3.744,15.45);
  \pic[rel] at (2.956,15.45) {un};
  \pic[rel] at (3.444,15.45) {patte};
\draw[rel] (8.056,15.45) -- (9.644,15.45);
  \pic[rel] at (8.206,15.45) {un};
  \pic[rel] at (9.344,15.45) {patte};
  \node[et] at (8.85,15.75) {0..N};
\draw[relg] (usb.south) -- (aph.north);
  \pic[relg, rotate=90]  at ($(usb.south)!0.30!(aph.north)$) {un};
  \pic[relg, rotate=-90] at ($(aph.north)!0.30!(usb.south)$) {patte};

% ================== bande de separation : aucune relation ====================
\node[note, text width=14.8cm] at (7.70,12.10)
  {\textbf{Aucune relation entre les deux cadres --- aucun arc ne franchit cette
   séparation.} L'identifiant d'analyste porté par les verdicts est l'identifiant
   technique du jeton, jamais l'adresse : le plan de supervision n'ouvre aucun accès
   à la base des comptes.};

% ================================ CADRE B ====================================
\draw[cadre] (0,5.00) rectangle (15.40,11.30);
\node[anchor=north west, font=\scriptsize\bfseries] at (0.15,11.24)
  {Cadre B --- entrepôt de supervision};
\node[anchor=north west, align=left, text width=15.10cm, font=\scriptsize] at (0.15,10.86)
  {\textbf{Aucune contrainte d'intégrité référentielle} : les relations ci-dessous sont des
   conventions de requête. \textbf{Six tables ne portent aucun arc} et ne figurent pas ici ---
   \texttt{raw\_logs}, \texttt{access\_logs}, \texttt{security\_events}, \texttt{api\_metrics}
   et \texttt{cve\_findings}, partitionnées par jour, et \texttt{pending\_embeddings},
   \emph{déclarée, jamais écrite ni lue} (détail : \cref{tab:modele-donnees}).};

% ---- coeur de la demonstration ----
\node[eh, text width=3.8cm] (deh) at (5.60,9.55) {detections};
\node[eb, text width=3.8cm] (deb) at (deh.south)
  {\texttt{id} \,\textperiodcentered\, horodatage \,\textperiodcentered\, \texttt{rule\_id}\\
   \texttt{\textbf{entity}}\\
   \texttt{service} = clé de locataire\\
   tactique \,\textperiodcentered\, technique};

\node[eh, text width=3.4cm] (emh) at (13.20,9.55) {attack\_embeddings};
\node[eb, text width=3.4cm] (emb) at (emh.south)
  {\texttt{technique} \textbf{PK}\\ vecteur à 768 dimensions};

\node[eh, text width=4.8cm] (enh) at (3.00,6.75) {alert\_enrichment};
\node[eb, text width=4.8cm] (enb) at (enh.south)
  {\texttt{detection\_id} \,\textperiodcentered\, technique \,\textperiodcentered\, tactique\\
   similarité \,\textperiodcentered\, version du modèle\\
   candidats de rang 2 et 3\\
   \textbf{aucune colonne \texttt{entity}}};

\node[eh, text width=4.4cm] (veh) at (12.40,6.75) {analyst\_verdicts};
\node[eb, text width=4.4cm] (veb) at (veh.south)
  {\texttt{\textbf{entity}}\\
   identifiant technique de l'analyste\\
   verdict \,\textperiodcentered\, \emph{ajout seul}};

% ---- traitement, hors cadre B ----
\node[trait, text width=5.8cm] (sco) at (6.90,4.30)
  {\textbf{Calcul du score d'incident}\\
   regroupement par \texttt{entity}, à la lecture};

% ---- relations du cadre B ----
% detections 1--N alert_enrichment
\draw[rel] (4.60,7.451) coordinate (q1) -- (3.60,6.75) coordinate (q2);
  \pic[rel, rotate=229] at ($(q1)!0.15cm!(q2)$) {un};
  \pic[rel, rotate=229] at ($(q2)!0.30cm!(q1)$) {patte};
  \node[etg, anchor=west] at (4.78,7.10) {par \texttt{detection\_id}};
% detections N--1 attack_embeddings, en tirete : rattachement sans cle
\draw[rel, densely dashed] (7.606,8.55) -- (11.394,8.55);
  \pic[rel, rotate=180] at (7.906,8.55) {patte};
  \pic[rel] at (11.244,8.55) {un};
  \node[et, text width=2.0cm] at (9.50,8.55)
    {rattachement par similarité,\\ non par clé};
% detections |><| analyst_verdicts, sur entity
\draw[rel] (7.606,8.00) -- (10.80,6.75);
  \node[et, text width=2.9cm] at (9.20,7.38)
    {$\bowtie$ \textbf{jointure du score}\\ sur \texttt{entity}};
% detections -> calcul du score
\draw[a] (7.15,7.451) -- (7.15,4.34);
  \node[etg, text width=2.7cm, anchor=west] at (7.30,5.45)
    {regroupement par \texttt{entity} : le score se calcule à la lecture};
% alert_enrichment ->< calcul du score : l'arc que le schema interdit
\draw[relg, densely dashed] (3.00,4.651) -- node[croixrefus]{} (3.00,3.60) -- (3.894,3.60);
  \node[etg, anchor=north west] at (0.10,3.40)
    {\textbf{jointure au score structurellement impossible} ---\\
     \texttt{alert\_enrichment} ne porte aucune colonne \texttt{entity}};

\end{tikzpicture}
\caption[Modèle de données des deux systèmes de stockage]{Modèle de données des deux systèmes de
stockage : ce que les clés permettent, et la jointure que leur absence interdit. L'impossibilité de
corréler le rattachement sémantique au score d'incident est une propriété \emph{du schéma} et non
un choix d'algorithme --- une colonne manquante suffit à l'expliquer --- et les deux systèmes de
stockage sont volontairement disjoints.}
\label{fig:modele-donnees}
\end{figure}

\section{De l'attaque à l'incident : fonctionnement de bout en bout}
\label{sec:bout-en-bout}

Entre la requête malveillante qui atteint le périmètre et l'incident qualifié que voit l'analyste,
deux registres se partagent le récit sans se recouvrir : la \cref{fig:seq-refus} dit \emph{qui}
agit et \emph{par quelle règle}, la \cref{fig:cycle-donnee} dit \emph{quand}, sa colonne
d'horodatage relatif situant chaque étape par rapport au refus opposé au périmètre.

La séquence porte une propriété du périmètre qu'aucune autre vue ne rend : l'évaluation de la
politique de filtrage est \textbf{premier-match-gagne} --- les règles sont ordonnées par priorité,
la première qui correspond décide, les suivantes ne sont pas évaluées. L'ordre est donc une
propriété de sécurité et non un détail de configuration : une règle de régulation de débit placée
avant les règles applicatives les rend inertes, puisqu'elle laisse passer de manière terminale le
trafic conforme au seuil --- c'est l'incident corrigé pendant la réalisation
(\cref{subsec:incidents}), et la raison pour laquelle la restriction géographique est écrite comme
un refus dont les exemptions d'administration sont repliées à l'intérieur.

\begin{figure}[htbp]
\centering
\begin{tikzpicture}[
  x=1cm, y=1cm, font=\scriptsize,
  part/.style={draw=black!70, fill=black!4, align=center,
               inner sep=2.5pt, minimum height=1.05cm},
  lif/.style={draw=black!70, densely dashed, line width=0.45pt},
  act/.style={draw=black!70, fill=black!10, line width=0.45pt},
  msg/.style={fluxdonnee},
  rep/.style={fluxreponse},
  asy/.style={fluxevenement},
  et/.style={font=\scriptsize, fill=white, inner sep=1pt, align=center},
  num/.style={circle, draw=black!70, fill=white, line width=0.4pt,
              inner sep=0.5pt, minimum size=3.4mm, font=\scriptsize},
  note/.style={draw=black!70, fill=black!4, align=left, inner sep=3.5pt,
               font=\scriptsize},
  lead/.style={draw=black!55, densely dotted, line width=0.4pt},
  brk/.style={draw=black!55, line width=0.4pt},
  cart/.style={draw=black!45, rounded corners=2pt, fill=black!2, align=left,
               inner sep=4pt, font=\scriptsize}]

% ===================== bandes de plan, calque de fond =========================
\begin{scope}[on background layer]
  \fill[menalplanA] (1.62,0.10) rectangle (5.20,-12.80);
  \fill[menalplanC] (5.30,0.10) rectangle (9.00,-12.80);
  \fill[menalplanB] (9.10,0.10) rectangle (15.02,-12.80);
\end{scope}

% ===================== participants ==========================================
\node[part, text width=1.00cm] (cl) at (2.30,-0.55) {Client};
\node[part, text width=1.45cm] (pe) at (4.28,-0.55) {Périmètre\\filtrage\\applicatif};
\node[part, text width=1.62cm] (jo) at (6.26,-0.55) {Journalisation};
\node[part, text width=1.25cm] (en) at (8.24,-0.55) {Entrepôt\\(journaux\\bruts)};
\node[part, text width=1.62cm] (no) at (10.22,-0.55) {Requête planifiée\\de normalisation};
\node[part, text width=1.55cm] (de) at (12.20,-0.55) {Requêtes planifiées\\de détection\\(R1--R7)};
\node[part, text width=1.62cm] (ap) at (14.05,-0.55) {API de supervision\\et tableau de bord};

% ===================== lignes de vie =========================================
\draw[lif] (2.30,-1.10)  -- (2.30,-12.55);
\draw[lif] (4.28,-1.10)  -- (4.28,-12.55);
\draw[lif] (6.26,-1.10)  -- (6.26,-12.55);
\draw[lif] (8.24,-1.10)  -- (8.24,-12.55);
\draw[lif] (10.22,-1.10) -- (10.22,-12.55);
\draw[lif] (12.20,-1.10) -- (12.20,-12.55);
\draw[lif] (14.05,-1.10) -- (14.05,-12.55);

% ===================== fragment ordonne (message 2) ==========================
% dessine avant les barres d'activation : son fond opaque interrompt les lignes
% de vie inactives de la Journalisation et de l'Entrepot sur cette plage.
\fill[white] (4.28,-2.25) rectangle (10.05,-6.20);
\draw[black!70, line width=0.5pt] (4.28,-2.25) rectangle (10.05,-6.20);
% etiquette du fragment, en coin superieur gauche (pentagone UML)
\fill[black!10] (4.28,-2.25) -- (7.45,-2.25) -- (7.45,-2.83)
                -- (7.27,-3.05) -- (4.28,-3.05) -- cycle;
\draw[black!70, line width=0.5pt] (4.28,-2.25) -- (7.45,-2.25) -- (7.45,-2.83)
                -- (7.27,-3.05) -- (4.28,-3.05) -- cycle;
\node[anchor=north west, align=left, inner sep=0pt] at (4.42,-2.34)
      {2 : évaluation ordonnée\\ \textbf{premier-match-gagne}};
% ordre d'evaluation, de haut en bas
\draw[msg] (4.60,-3.28) -- (4.60,-4.28);
\node[anchor=west, inner sep=0pt] at (4.82,-3.38) {410 \ restriction géographique};
\node[anchor=west, inner sep=0pt] at (4.82,-3.73) {1000 \ XSS};
\node[anchor=west, inner sep=0pt] at (4.82,-4.08)
      {1100 \ injection SQL \ \
       \tikz[baseline=-0.55ex]{\draw[black!70, line width=0.7pt]
         (0,0.05) -- (0.07,-0.04) -- (0.19,0.14);}\;\textbf{correspond}};
\draw[black!70, line width=0.4pt] (4.28,-4.38) -- (10.05,-4.38);
\node[anchor=north west, text width=5.35cm, align=left, inner sep=0pt,
      text=menalencre] at (4.48,-4.54)
      {règles suivantes \textbf{non évaluées} : 1200 LFI, 1300 RCE, 1400 RFI,
       1450 débit d'authentification, 1500 régulation globale, autorisation
       par défaut};

% ---- note rattachee au fragment (trait fin) ---------------------------------
\node[note, text width=5.15cm, anchor=north west] (nt) at (10.35,-2.35)
      {une règle de régulation de débit placée \emph{avant} les règles
       applicatives est terminale et les rend inertes~; c'est la raison du
       placement en priorité 1500.};
\fill[white] (nt.north east) -- ++(-0.30,0) -- ++(0,-0.30) -- cycle;
\draw[black!70, line width=0.5pt]
      ([xshift=-0.30cm]nt.north east) -- ([yshift=-0.30cm]nt.north east);
\draw[lead] (10.05,-3.10) -- (nt.west);

% ===================== barres d'activation ===================================
\draw[act] (2.19,-1.65)   rectangle (2.41,-6.65);
\draw[act] (4.17,-1.72)   rectangle (4.39,-7.45);
\draw[act] (6.15,-7.25)   rectangle (6.37,-8.20);
\draw[act] (8.13,-7.90)   rectangle (8.35,-8.25);
\draw[act] (8.13,-8.85)   rectangle (8.35,-9.20);
\draw[act] (8.13,-9.90)   rectangle (8.35,-10.25);
\draw[act] (8.13,-10.95)  rectangle (8.35,-11.30);
\draw[act] (10.11,-8.85)  rectangle (10.33,-9.20);
\draw[act] (12.09,-9.90)  rectangle (12.31,-10.25);
\draw[act] (13.94,-10.95) rectangle (14.16,-12.15);

% ===================== messages ==============================================
% 1 -- requete malveillante
\draw[msg] (2.41,-1.70) -- (4.17,-1.70);
\node[num] at (2.72,-1.70) {1};
\node[et] at (3.75,-1.34) {requête HTTPS malveillante (F1)};

% 3 -- reponse 403 : c'est un REFUS, il porte le trait et la croix de refus
\draw[fluxrefuse] (4.17,-6.55) -- node[croixrefus]{} (2.41,-6.55);
\node[num] at (3.86,-6.55) {3};
\node[et] at (3.20,-6.79) {\textbf{403}};

% 4 -- verdict vers la journalisation (asynchrone)
\draw[asy] (4.39,-7.30) -- (6.15,-7.30);
\node[num] at (4.70,-7.30) {4};
\node[et] at (6.05,-6.94) {verdict, échantillonnage intégral};

% 5 -- journalisation vers entrepot (asynchrone)
\draw[asy] (6.37,-8.00) -- (8.13,-8.00);
\node[num] at (6.68,-8.00) {5};

% 6 -- normalisation
\draw[msg] (10.11,-9.00) -- (8.35,-9.00);
\node[num] at (9.80,-9.00) {6};
\node[et] at (9.05,-8.64) {fenêtre glissante 15~min};
\node[et] at (9.15,-9.38) {\emph{cadence 5~min}};

% 7 -- detection
\draw[msg] (12.09,-10.05) -- (8.35,-10.05);
\node[num] at (11.64,-10.05) {7};
\node[et] at (10.10,-9.69) {R1--R7, insertion conditionnelle};
\node[et] at (10.10,-10.43) {\emph{cadence 5~min, fenêtre 15~min}};

% 8 -- lecture par l'API
\draw[msg] (13.94,-11.10) -- (8.35,-11.10);
\node[num] at (13.27,-11.10) {8};
\node[et] at (11.00,-10.74) {regroupement par entité, score calculé à la lecture};

% 9 -- restitution a l'analyste : message vers l'exterieur du diagramme
\draw[msg] (14.16,-12.00) -- (15.30,-12.00);
\fill[black!70] (15.33,-12.00) circle (1.3pt);
\node[num] at (14.47,-12.00) {9};
\node[et] at (11.90,-11.64) {incident affiché --- vers l'analyste};
\node[et] at (12.60,-12.38) {\emph{rafraîchissement 10~s}};

% ===================== colonne des delais ====================================
\draw[brk] (1.98,-1.65) -- (1.80,-1.65) -- (1.80,-6.65) -- (1.98,-6.65);
\node[text width=1.90cm, align=left, anchor=east, inner sep=0pt] at (1.72,-4.15)
      {\textbf{1 à 3}\\ synchrone, fraction de seconde --- \emph{configuration,
       non une latence mesurée}};
\draw[brk] (1.98,-7.25) -- (1.80,-7.25) -- (1.80,-8.05) -- (1.98,-8.05);
\node[text width=1.90cm, align=left, anchor=east, inner sep=0pt] at (1.72,-7.65)
      {\textbf{4 et 5}\\ quelques secondes --- \emph{configuration, non
       mesurée}};

% ===================== notation ==============================================
\node[anchor=north west, inner sep=0pt, align=left, text width=15.40cm]
      at (-0.18,-12.95)
  {\textbf{Notation}\;
   \tikz[baseline=-0.55ex]{\draw[msg] (0,0) -- (0.62,0);}\,appel synchrone \quad
   \tikz[baseline=-0.55ex]{\draw[rep] (0,0) -- (0.62,0);}\,réponse \quad
   \legendecroix{refus} \quad
   \tikz[baseline=-0.55ex]{\draw[asy] (0,0) -- (0.62,0);}\,message asynchrone \quad
   \tikz[baseline=-0.55ex]{\draw[lead] (0,0) -- (0.62,0);}\,note \quad
   \emph{italique} : cadence configurée \quad
   fonds, de gauche à droite : plan de données, plan d'observation, plan de contrôle};

\end{tikzpicture}
\caption[Requête malveillante refusée au périmètre]{Requête malveillante refusée au périmètre :
ordre d'évaluation des règles et trajet de la preuve jusqu'à l'analyste. L'ordre est une propriété
de sécurité --- la première correspondance est terminale et les six règles suivantes ne sont jamais
évaluées --- et la preuve du refus met beaucoup plus de temps à atteindre l'analyste que le refus
lui-même à être prononcé.}
\label{fig:seq-refus}
\end{figure}


Deux définitions commandent leur lecture : une \textbf{requête planifiée} est une requête
enregistrée dans l'entrepôt et exécutée à intervalle fixe sans intervention humaine, et la
\textbf{fenêtre glissante} est l'intervalle immédiatement passé sur lequel porte chaque exécution.
S'y ajoute le \textbf{démarrage à froid}, mesuré au \cref{subsec:incidents}, dont dépend avec la
cadence de quinze minutes le délai d'enrichissement (F5).

\textbf{Deux mesures, et deux seulement, encadrent le délai de qualification.} Observation du
19/08/2026 : 13 requêtes bloquées en 403 à 16~h~41 UTC, règle R2 déclenchée à 16~h~56~min~08~s UTC,
comptage exact de 13 --- soit environ 15~min. Observation du 10/08/2026 : 6~min~15~s sur un
scénario de force brute (R1). \textbf{Aucune des deux n'est un maximum garanti} ; toutes les autres
valeurs portées par les figures décrivent la configuration, non une latence mesurée.

\section{Conception détaillée par couche}
\label{sec:conception-couches}

Le \cref{tab:conception-couches-synthese} donne les sept niveaux, les composants qui les
constituent et les exigences qu'ils couvrent --- les capacités restant le vocabulaire du rapport,
les produits n'en étant que l'instanciation chez le fournisseur retenu par D00. Le
\textbf{détail des contrôles apportés par chaque niveau} est publié en
\cref{ann:conception-couches}, dans le même ordre. Le nombre d'exigences ne mesure pas la
criticité d'une couche : L5 est la moins chargée et porte pourtant toute la chaîne de preuve, des
journaux au score.

\begingroup
\setlength{\tabcolsep}{3pt}
\renewcommand{\arraystretch}{1.0}
\begin{longtable}{|P{2.0cm}|P{11.0cm}|P{1.9cm}|}
\caption[Les sept niveaux : composants et exigences couvertes]{Les sept niveaux de la conception :
composants employés, produits qui les réalisent et exigences couvertes. Les contrôles apportés par
chaque niveau sont détaillés en \cref{ann:conception-couches}.}
\label{tab:conception-couches-synthese} \\
\hline
\textbf{Niveau} & \textbf{Composants et produits employés} & \textbf{Exigences} \\
\hline
\endfirsthead
\hline
\textbf{Niveau} & \textbf{Composants et produits employés} & \textbf{Exigences} \\
\hline
\endhead
\textbf{L1}\newline Périmètre d'entrée &
Répartiteur global HTTPS, trois certificats gérés, moteur de filtrage applicatif, limitation de
débit, restriction géographique. \emph{Cloud Load Balancing}, \emph{Cloud Armor} --- politique
unique attachée à tous les services publiés. La zone de noms est une dépendance externe (É10) &
EX1, EX2, EX3 \\ \hline
\textbf{L2}\newline Identité et sécurité \emph{(plan transversal)} &
\emph{Cloud IAM} (sept identités de service), \emph{Secret Manager} (neuf secrets, dont huit
chiffrés par une clé gérée par le projet), \emph{Cloud KMS} (deux trousseaux, deux clés, rotation
de quatre-vingt-dix jours), \emph{Workload Identity Federation} &
EX5, EX6, EX7, EX10, EX11 \\ \hline
\textbf{L3}\newline Charges de travail &
\textbf{Cinq services et trois tâches} : API de supervision, tableau de bord, service d'encodage,
interface web et interface de programmation de l'application hébergée ; enrichissement, migration
du schéma, vérification quotidienne de l'isolation des bases. \emph{Cloud Run}, \emph{Cloud
Scheduler} pour deux déclencheurs, la chaîne de livraison pour le troisième &
EX6, EX9 \\ \hline
\textbf{L4}\newline Réseau &
Un réseau virtuel en mode personnalisé, \textbf{un sous-réseau par locataire} dont un réservé au
troisième, un appairage de services privés vers la base, et \textbf{douze règles de filtrage dont
cinq autorisations}. \emph{VPC}, \emph{sortie réseau directe} et \emph{Private Service Access} &
EX4, EX20 \\ \hline
\textbf{L5}\newline Données &
\textbf{Base applicative} (comptes, rôles, journal d'audit) et \textbf{entrepôt de supervision},
c\oe ur du dispositif de détection (\cref{fig:cycle-donnee}, \cref{fig:modele-donnees}).
\emph{Cloud SQL for PostgreSQL 15} (adresse privée seule, haute disponibilité régionale,
restauration à un instant donné sur sept jours), \emph{BigQuery} (dix tables), \emph{Cloud
Storage} (clé gérée, prévention d'accès public) &
EX2, EX4, EX11 \\ \hline
\textbf{L6}\newline Enrichissement sémantique &
Un service d'encodage sans état, une tâche d'enrichissement par lots, un référentiel de techniques
d'attaque vectorisé. \emph{Cloud Run} en entrée interne seule, modèle de représentation de phrases
spécialisé en cybersécurité embarqué dans l'image et exécuté hors ligne, recherche vectorielle
native de l'entrepôt --- ce qui supprime la base vectorielle séparée qu'imposerait un autre
fournisseur (D05, D06) &
EX14, EX15 \\ \hline
\textbf{L7}\newline Observabilité \emph{(plan transversal)} &
Quatre exports de journaux, politiques d'alerte, objectifs de service, sondes de disponibilité,
lien entre une alerte affichée et sa procédure d'incident. \emph{Cloud Logging}, \emph{Cloud
Monitoring} &
EX12, EX13, EX19 \\
\hline
\end{longtable}
\endgroup

\textbf{Points de vigilance.} Chaque niveau nomme ce qu'il ne garantit pas lui-même, et quatre des
sept --- L1, L3, L4 et L6 --- y désignent une \textbf{confiance résiduelle} : un contrôle qu'ils
n'exercent pas et qu'un contrôle voisin doit reprendre. Les deux contrôles dynamiques qui les
établissent en propre, T14 et T20, portent leur relevé daté et figurent au \textbf{plan de
validation continue} (\cref{sec:plan-validation}).

\textbf{L1} charge les règles préconfigurées à leur niveau de sensibilité le plus bas : arbitrage
assumé entre un taux de faux positifs intenable pour un exploitant unique et la résistance aux
techniques d'évasion, mesurée au \cref{ch:validation}.

\textbf{L2} ne peut interdire structurellement la création d'une clé d'identité de longue durée ---
cela relève d'une politique d'organisation, dont la pose exige un périmètre administratif dont le
projet ne dispose pas (É7) : le contrôle est porté par la fédération, qui rend la clé inutile, et
vérifié par inventaire. Seconde réserve, de conception : la clé qui chiffre les secrets du second
facteur ne porte pas d'identifiant de version, de sorte qu'une rotation les rendrait illisibles d'un
seul coup ; son versionnement est la première évolution prévue de ce plan.

\textbf{L3} accorde le droit d'invocation à tous sur les services publiés : leur protection tient
entièrement au paramètre d'entrée restreinte au répartiteur, contrôle unique dont une régression les
exposerait publiquement. Le choix est assumé dans le code, faute d'un service d'authentification en
amont disponible hors organisation cloud ; le contrôle compensatoire est la vérification de ce champ
à chaque plan d'infrastructure, le seuil de bascule l'acquisition d'une organisation. La limite de
requêtes simultanées par instance n'est déclarée nulle part : seule valeur du modèle d'exécution qui
échappe à PR4.

\textbf{L4} ne porte plus la réserve qu'enregistrait D11 : depuis la refonte appliquée du 11 au
25/08/2026 (\cref{subsec:reseau-cible}), la sortie est gouvernée par le pare-feu et non par un
paramètre de plateforme, la règle interne de priorité 900 a été supprimée le 24/08/2026 avec les
deux sous-réseaux vides qu'elle reliait (É2), et le canal de détection des refus réseau n'est plus
muet --- la journalisation des refus, acquise le 19/08/2026 (É1), a porté le 24/08/2026 trois refus
de sortie réellement opposés et retrouvés dans l'entrepôt (T20). \textbf{La confiance résiduelle de
la couche a changé de nature, et elle est mesurée} : le service d'encodage et l'interface web de
l'application hébergée ont été détachés du sous-réseau au motif, exact, qu'ils n'ont aucune
dépendance privée, de sorte que leur \textbf{sortie managée échappe au pare-feu}. Le 25/08/2026, une
sonde portant la configuration réseau de la révision déployée du service d'encodage a atteint une
destination externe en 412~ms : c'est la non-conformité T14, qui reste ouverte, le motif du
détachement valant pour la sortie \emph{nécessaire} et non pour la sortie \emph{possible}.

\textbf{L5} n'applique pas rétroactivement le chiffrement par clé gérée : deux tables seulement
portent une configuration explicite, les tables de preuves antérieures restant chiffrées par la clé
du fournisseur, et généraliser supposerait de les recréer donc d'interrompre la chaîne de preuve ---
arbitrage tranché en faveur de la continuité. Point de méthode : abaisser le seuil de rattachement
jusqu'à ce que toute alerte en reçoive un aurait produit un tableau de bord entièrement rempli et
entièrement faux ; le choix inverse a été fait, afficher le taux d'alertes non rattachées comme
indicateur de qualité.

\textbf{L6} porte au premier chef la sortie non gouvernée du service d'encodage, énoncée en L4.
Second point, dimensionnant : la recherche vectorielle s'exécute sans index déclaré, chaque
appariement balayant l'intégralité du référentiel --- acceptable à la volumétrie d'ATT\&CK et à la
cadence de quinze minutes, c'est la première limite qu'une montée en charge rencontrerait.

\textbf{L7} laisse sans filtre le quatrième export, poste de coût dominant de l'entrepôt, parce que
c'est lui qui rend la détection possible : les verdicts du filtrage applicatif n'existent que dans ce
journal, et tout échantillonnage y créerait un angle mort statistique sur la source de preuve la plus
dense. La condition de réévaluation est le franchissement du budget d'ingestion mensuel.

\subsection{Couverture des exigences de sécurité}
\label{sec:couverture-exigences}

Les vingt exigences du \cref{ch:besoins-menaces} sont toutes prises en charge par au moins une couche ou un flux,
le détail figurant en annexe~\ref{ann:conception-detaillee} ; cette prise en charge n'est pas une
preuve de conformité, implémentation et niveau de test étant évalués au \cref{ch:validation}. Trois réserves,
portant sur quatre exigences, sont nommées : \textbf{EX7}, politique d'organisation inapplicable
faute d'organisation cloud (É7), compensée par les autorisations et l'infrastructure en code ;
\textbf{EX9}, déploiement par étiquette égale au SHA du commit et non par empreinte d'image stricte
(É8) ; \textbf{EX16 et EX18}, portée bornée aux ressources déjà décrites en code et revue
d'approbation manuelle. La correspondance avec les sept principes Zero Trust du \cref{ch:etat-art} a été
vérifiée point par point : l'identité comme plan de contrôle et la mesure continue de l'état y sont
les mieux représentés.

\section{Le plan d'identité : la segmentation réelle}
\label{sec:plan-identite}

Cette section traduit le principe d'identité (PR2) : ce qui, dans cette architecture, remplace les
pare-feux internes.

\subsection{Matrice des identités et des droits}

Le \cref{tab:matrice-identites} recense les identités du socle et leurs droits ; la
\cref{fig:plan-identite} en projette les cinq porteuses d'un droit direct sur une ressource,
et rend visible ce qu'un tableau montre mal : les interdictions.

\begin{table}[htbp]
\caption{Matrice des identités et des droits.}
\label{tab:matrice-identites}
\centering
\footnotesize
\setlength{\tabcolsep}{3.5pt}
\renewcommand{\arraystretch}{1.05}
\begin{tabular}{|P{3.4cm}|P{6.1cm}|P{5.6cm}|}
\hline
Identité & Droits accordés & Interdiction structurante \\
\hline
Identité applicative (une par application) & Accès à la base de données ; accès aux secrets de
cette application uniquement & \textbf{Aucun accès à l'entrepôt de supervision} \\ \hline
Identité de l'API de supervision & Lecture seule sur l'entrepôt ; accès à la base ; secrets de
la plateforme & Aucune écriture sur les preuves --- exception unique : la table des verdicts
d'analyste \\ \hline
Identité du tableau de bord & Appel de l'API de supervision ; lecture du jeu de données de
supervision & \textbf{Aucun accès à la base de données applicative} ; séparation vis-à-vis de
l'entrepôt appliquée par le code et non par les autorisations (écart É9) \\ \hline
Identité de la tâche d'enrichissement & Lecture des détections ; écriture \textbf{uniquement}
sur la table d'enrichissement ; appel du service d'encodage & \textbf{Écriture interdite sur les
journaux bruts et les détections} \\ \hline
Identité de déploiement & Écriture dans le registre ; déploiement ; usurpation
contrôlée des identités d'exécution & Aucune clé exportée ; portée limitée au projet de son
environnement \\
\hline
\end{tabular}
\end{table}

La plus significative est celle de la tâche d'enrichissement : l'interdiction d'écriture sur les
détections empêche un attaquant qui l'aurait compromise d'effacer la trace de son passage (scénario
SO16, test T11 du \cref{ch:validation}).

\begin{figure}[htbp]
\centering
\begin{tikzpicture}[
  iden/.style={draw, thick, rounded corners=3pt, fill=black!4, text width=3.1cm, align=center,
    font=\footnotesize, inner sep=3pt, minimum height=0.75cm},
  res/.style={draw, thick, fill=white, text width=3.1cm, align=center,
    font=\footnotesize, inner sep=3pt, minimum height=0.75cm},
  ok/.style={fluxdonnee},
  ko/.style={fluxrefuse},
  e9/.style={lienstructure, densely dashed, -{Latex[length=1.6mm]}},
  etqa/.style={font=\scriptsize, above=1pt, fill=white, inner sep=1pt}
]
\node[font=\footnotesize\bfseries] at (0,6.3)   {Identités};
\node[font=\footnotesize\bfseries] at (8.4,6.3) {Ressources};
\node[iden] (e) at (0,5.40) {Tableau de bord};
\node[iden] (a) at (0,4.05) {Application hébergée};
\node[iden] (b) at (0,2.70) {API de supervision};
\node[iden] (c) at (0,1.35) {Tâche\\d'enrichissement};
\node[iden] (d) at (0,0.00) {Identité de déploiement};
\node[res] (r5) at (8.4,5.40) {API de supervision};
\node[res] (r1) at (8.4,4.05) {Base de données};
\node[res] (r2) at (8.4,2.70) {Journaux bruts et détections};
\node[res] (r3) at (8.4,1.35) {Table d'enrichissement};
\node[res] (r4) at (8.4,0.00) {Registre d'images};
% Ordre de trace volontaire : les diagonales d'abord, les liaisons horizontales
% et leurs libelles ensuite. TikZ dessine dans l'ordre du source ; placees apres,
% les diagonales rayaient le fond blanc des libelles deja poses.
\draw[ko] (a) -- node[croixrefus]{} (r2);
\draw[ko] (c) -- node[croixrefus]{} (r2);
\draw[ko] (e) -- node[croixrefus]{} (r1);
% pos 0.82 plutot que 0.70 : a 0.70 l'etiquette tombait exactement sur le
% croisement de cette diagonale avec la liaison « lecture / ecriture ».
% text=black : le trait reste gris, le libelle doit rester lisible.
\draw[e9] (e) -- (r2)
  node[etqa, text=black, pos=0.82] {non cloisonné (É9)};
\draw[ok] (a) -- (r1) node[etqa, pos=0.28] {lecture / écriture};
\draw[ok] (b) -- (r2) node[etqa, pos=0.3] {lecture seule};
\draw[ok] (c) -- (r3) node[etqa, pos=0.3] {écriture};
\draw[ok] (d) -- (r4) node[etqa, pos=0.3] {écriture};
\draw[ok] (e) -- (r5) node[etqa, pos=0.3] {appel API};
% --- separation des plans : les identites sont le plan de controle, les
%     ressources le plan de donnees. Calque de fond, aucun surcout de hauteur.
\begin{scope}[on background layer]
  \node[bandecontrole, fit=(e)(a)(b)(c)(d)]           {};
  \node[bandedonnee,   fit=(r5)(r1)(r2)(r3)(r4)]      {};
\end{scope}
\node[legendecadre, anchor=north west, font=\scriptsize, text width=12.0cm] at (-1.7,-0.75)
  {\legendetrait{fluxdonnee}{droit accordé}\quad
   \legendecroix{interdiction structurante}\quad
   \legendetrait{e9}{séparation appliquée par le code (É9)}\\[2pt]
   Fond des colonnes : à gauche le plan de contrôle (les identités), à droite le plan de données};
\end{tikzpicture}
\caption[Plan d'identité : droits accordés et interdictions structurantes]{Plan d'identité : droits
accordés et interdictions structurantes (cinq identités). Ni l'application hébergée ni la tâche
d'enrichissement ne peuvent écrire sur les preuves (scénario SO16, test T11) ; le tableau de bord
passe par l'API de supervision, mais son identité conserve un droit de lecture sur l'entrepôt ---
c'est l'écart É9.}%
\label{fig:plan-identite}
\end{figure}

\subsection{Chaînes de délégation d'identité}
\label{subsec:delegation}

Une matrice de droits dit ce que chaque identité peut faire, non ce qu'un attaquant obtient en
\emph{enchaînant} les délégations : c'est l'objet de la \cref{fig:delegation-identite}, dont il faut
regarder le \emph{dernier} segment. \textbf{L'identité de déploiement porte un droit de déploiement
à l'échelle du projet, et non ressource par ressource} : combinée à l'usurpation de l'identité de
l'API de supervision --- qui accède à la clé de signature des jetons de session --- elle ouvre un
chemin allant de la branche principale du dépôt à un jeton portant le rôle d'administration. Le code
a fermé la porte adjacente, en refusant à cette identité d'API tout accès aux secrets à l'échelle du
projet et en le limitant à trois secrets nommés ; la chaîne subsiste par le droit de déploiement. La
réduction retenue --- restreindre ce droit aux seuls services réellement déployés par la chaîne,
ressource par ressource --- est portée au plan de remédiation et conditionne l'accueil d'un
troisième locataire. Ce chemin est nommé ici parce qu'un modèle de menace qui ne considère que les
droits directs manque systématiquement les chemins composés.

\begin{figure}[htbp]
\centering
% --- pastilles dessinees en TikZ : aucune dependance a un glyphe Unicode ------
\providecommand{\pastilleOK}{\tikz[baseline=-0.55ex]{%
  \draw[black!70, line width=0.30pt] (0,0) circle (0.130cm);
  \draw[black!70, line width=0.60pt, line cap=round, line join=miter]
       (-0.055cm,-0.005cm) -- (-0.020cm,-0.072cm) -- (0.085cm,0.095cm);}}
\providecommand{\pastilleWarn}{\tikz[baseline=-0.55ex]{%
  \draw[black!70, line width=0.40pt, line join=miter]
       (0,0.125cm) -- (0.132cm,-0.105cm) -- (-0.132cm,-0.105cm) -- cycle;
  \draw[black!70, line width=0.55pt, line cap=butt] (0,0.058cm) -- (0,-0.018cm);
  \fill[black!70] (0,-0.058cm) circle (0.021cm);}}

\begin{tikzpicture}[
  x=1cm, y=1cm, font=\scriptsize,
  etape/.style={draw=black!70, rounded corners=3pt, fill=black!4,
                align=left, inner sep=3pt},
  ident/.style={draw=black!70, fill=white, align=left, inner sep=3pt},
  res/.style={draw=black!70, double, double distance=0.7pt, fill=white,
              align=left, inner sep=3pt},
  sec/.style={draw=black!55, fill=white, align=left, inner sep=2.5pt},
  encart/.style={draw=black!45, densely dotted, rounded corners=2pt, fill=black!2,
                 align=left, inner sep=3.5pt},
  num/.style={draw=black!70, circle, inner sep=1pt, font=\scriptsize\bfseries},
  a/.style={fluxdonnee},
  af/.style={fluxcontrole},
  trunk/.style={lienstructure, line width=0.9pt},
  lead/.style={black!55, densely dotted},
  synth/.style={fluxsynthese},
  et/.style={font=\scriptsize, fill=white, inner sep=1.5pt}]

% =============================================================================
% COLONNES 1 A 5 -- lecture de gauche a droite
% =============================================================================
\node[etape, text width=2.30cm, minimum width=2.55cm] (c1) at (1.275,11.50)
  {\textbf{1. Forge logicielle}\\[1.5pt]
   exécution déclenchée par une poussée sur \texttt{main}};

\node[etape, text width=2.30cm, minimum width=2.55cm] (c2) at (4.575,11.50)
  {\textbf{2. Fédération d'identité}\\[1.5pt]
   assertion OIDC\\[1.5pt]
   \emph{condition d'attribut :}\\
   \textbf{dépôt ET branche}};

\node[etape, text width=2.55cm, minimum width=2.80cm] (c3) at (8.000,11.50)
  {\textbf{3. Identité de déploiement}\\[2pt]
   \pastilleWarn~droit de déploiement \textbf{à l'échelle du projet}\\[2pt]
   écriture au registre à l'échelle du projet\\[2pt]
   \pastilleOK~écriture sur la \textbf{seule} table des vulnérabilités};

\node[ident, text width=1.55cm, minimum width=1.80cm] (i1) at (11.40,13.20)
  {\textbf{4a.} Identité de l'API de supervision};
\node[ident, text width=1.55cm, minimum width=1.80cm] (i2) at (11.40,11.50)
  {\textbf{4b.} Identité du tableau de bord};
\node[ident, text width=1.55cm, minimum width=1.80cm] (i3) at (11.40,9.80)
  {\textbf{4c.} Identité de l'application hébergée};

\node[res, text width=1.95cm, minimum width=2.20cm] (r1) at (14.50,14.60)
  {Clé de signature des jetons de session};
\node[res, text width=1.95cm, minimum width=2.20cm] (r2) at (14.50,13.05)
  {Clé de chiffrement du secret du second facteur};
\node[res, text width=1.95cm, minimum width=2.20cm] (r3) at (14.50,11.50)
  {Mot de passe de la base};
\node[res, text width=1.95cm, minimum width=2.20cm] (r4) at (14.50,9.95)
  {Registre d'images};
\node[res, text width=1.95cm, minimum width=2.20cm] (r5) at (14.50,8.40)
  {Médias de l'application hébergée};

% =============================================================================
% ARCS DU GRAPHE PRINCIPAL
% =============================================================================
\draw[a] (c1.east) -- (c2.west);
\node[et, rotate=90, align=center] at (2.925,12.90) {assertion\\signée};

\draw[a] (c2.east) -- (c3.west);
\node[et, rotate=90, align=center] at (6.225,12.90) {jeton de\\courte durée};

% -- 3 vers 4 : tronc vertical, trois sorties horizontales
\draw[trunk] (9.95,9.80) -- (9.95,13.20);
\draw[af]    (c3.east) -- (9.95,11.50);
\draw[a] (9.95,13.20) -- (i1.west);
\draw[a] (9.95,11.50) -- (i2.west);
\draw[a] (9.95,9.80)  -- (i3.west);
\node[align=center, text width=2.40cm, font=\scriptsize] (lus) at (9.95,14.38)
  {usurpation autorisée ($\times$3)};
\draw[lead] (9.95,14.02) -- (9.95,13.20);

% -- 4 vers 5 : tronc vertical, cinq sorties horizontales
\draw[trunk] (12.85,8.40) -- (12.85,14.60);
\draw[af] (i1.east) -- (12.85,13.20);
\draw[af] (i2.east) -- (12.85,11.50);
\draw[af] (i3.east) -- (12.85,9.80);
\draw[a] (12.85,14.60) -- (r1.west);
\draw[a] (12.85,13.05) -- (r2.west);
\draw[a] (12.85,11.50) -- (r3.west);
\draw[a] (12.85,9.95)  -- (r4.west);
\draw[a] (12.85,8.40)  -- (r5.west);
\node[align=center, text width=4.40cm, font=\scriptsize, anchor=north] (las) at (12.85,7.80)
  {accès à trois secrets \emph{nommés}\\--- \textbf{jamais} un droit à l'échelle du projet};
\draw[lead] (12.85,7.80) -- (12.85,8.40);

% =============================================================================
% ARC DE SYNTHESE -- en surimpression, au-dessus du graphe
% =============================================================================
\draw[synth] (c1.north) -- (1.275,16.40) -- (14.50,16.40) -- (r1.north);
\node[anchor=south, align=center, text width=11.0cm, font=\scriptsize, fill=white,
      inner sep=2pt] at (7.80,16.52)
  {\textbf{chemin composé} : pousser sur \texttt{main} $\rightarrow$ déployer une révision
   arbitraire $\rightarrow$ lire la clé $\rightarrow$ forger un jeton de rôle d'administration};

% =============================================================================
% ENCARTS
% =============================================================================
\node[encart, text width=5.35cm, anchor=north] (e2) at (2.800,9.30)
  {\pastilleOK~aucune clé de longue durée, vérifiée par inventaire le 19/08/2026\\[2.5pt]
   Sans la contrainte de branche, toute branche obtiendrait l'identité : le garde-fou
   posé côté forge n'est pas opposable au fournisseur.};
\draw[lead] (c2.south) -- (3.60,9.30);

\node[encart, text width=3.95cm, anchor=north] (e3) at (8.150,9.30)
  {\textbf{Réduction} --- restreindre le droit de déploiement aux seuls services
   réellement déployés, ressource par ressource ; condition : accueil d'un
   troisième locataire.};
\draw[lead] (c3.south) -- (8.150,9.30);

% =============================================================================
% LEGENDE INTERNE
% =============================================================================
\node[legendecadre, anchor=north west, font=\scriptsize, text width=15.4cm] at (0,7.10)
  {\legendetrait{fluxdonnee}{étape de la chaîne}\quad
   \legendetrait{fluxcontrole}{atteinte d'une ressource par un droit délégué}\quad
   \legendetrait{lienstructure}{rattachement}\\[2pt]
   \legendetrait{fluxsynthese}{le chemin composé --- accent unique de la figure}};

\end{tikzpicture}
\caption[Chaînes de délégation d'identité]{Chaînes de délégation d'identité : ce qu'un droit de
déploiement à l'échelle du projet rend atteignable. Aucun droit pris isolément n'est anormal --- la
fédération est durcie par une condition d'attribut portant sur le dépôt \emph{et} sur la branche,
sans clé de longue durée, et l'écriture à l'entrepôt est bornée à une seule table. Le risque naît
de la \textbf{composition} : c'est la portée du droit de déploiement, en aval, qui porte le risque
résiduel, non la fédération, en amont.}
\label{fig:delegation-identite}
\end{figure}

\textbf{Trois chaînes secondaires de délégation} suivent le même motif --- un agent managé du
fournisseur obtient le droit de créateur de jeton sur une identité d'exécution : l'agent de
transfert des requêtes planifiées et celui de l'ordonnanceur vers l'identité des règles de
détection, et l'opérateur humain vers l'identité de la tâche d'enrichissement lors du rejeu des
tests de bout en bout. Cette dernière liaison est nominative et \textbf{posée hors du code} :
c'est la seule des trois qu'aucun plan d'infrastructure ne rejouerait.

\subsection{Fédération de la chaîne de livraison, secrets et chiffrement}

Le point bloquant B2 concernait un secret unique gouvernant sept mécanismes ; le scénario SO9, une
clé d'identité de longue durée détenue par la chaîne de livraison. La réponse aux deux est de même
nature : \textbf{supprimer la classe de risque plutôt que la gérer.} La chaîne de livraison
s'authentifie par fédération d'identité --- une assertion signée par la forge, vérifiée par le
fournisseur cloud, contre un jeton de courte durée --- sous une condition d'attribut restreignant la
fédération à un dépôt et une branche précis (\cref{subsec:wif}). La politique interdisant la
création de clés de compte de service étant inapplicable faute d'organisation cloud (É7), l'absence
de clé est vérifiée par inventaire, non garantie par une contrainte structurelle. Deux règles de
secrets complètent ce plan, sans exception : \textbf{un secret par usage, jamais de secret dérivé
d'un autre} --- B2 en démontrait le coût, la fuite d'une valeur unique compromettant
l'authentification, l'anonymat et l'actif principal --- et \textbf{l'infrastructure en code crée les
conteneurs de secrets, jamais leurs valeurs}, injectées hors du code, pour que le fichier d'état ne
devienne jamais un coffre.

\subsection{Authentification à deux facteurs et cycle de la session}

L'accès humain au tableau de bord est le seul point du socle où une personne, et non une charge de
travail, présente une identité. Ce qu'il faut regarder dans la \cref{fig:seq-mfa} est le
\textbf{jeton intermédiaire} : entre la vérification du mot de passe et celle du code à usage
unique, la session est authentifiée sans être autorisée. Ce jeton ne porte \emph{aucune revendication
de rôle}, vit cinq minutes et est refusé par toutes les routes métier ; symétriquement, un jeton
d'accès est refusé par la route de vérification du second facteur, et le rôle n'apparaît qu'au
cinquième état, \emph{relu en base} plutôt que recopié. C'est cette séparation des types de jeton,
plus que le second facteur, qui constitue le contrôle.

Deux conséquences de conception en découlent. \textbf{La limitation de débit est à deux niveaux} :
au périmètre, dix requêtes par 60~s et par adresse avec bannissement de 300~s ; en applicatif, un
compteur indexé sur le \emph{défi} en cours et non sur l'adresse source, parce que l'appel relayé
présente l'adresse de sortie partagée du tableau de bord. Les deux couches ne mesurent donc pas la
même chose, et le raisonnement qui conduit à ce choix, mesure à l'appui, est exposé au
\cref{subsec:authent-dashboard}. La session est par ailleurs un jeton signé sans état serveur,
donc non révocable avant son échéance d'une heure : arbitrage assumé --- l'état serveur
qu'exigerait la révocation ajouterait un composant à exploiter, ce que PR1 écarte à ce
périmètre --- dont la condition de réévaluation est l'ouverture du tableau de bord à des comptes
autres que ceux de l'exploitant unique.

\begin{figure}[htbp]
\centering
\begin{tikzpicture}[
  x=1cm, y=1cm, font=\scriptsize,
  part/.style={draw=black!70, fill=black!4, align=center,
               inner sep=2.5pt, minimum height=0.95cm},
  lif/.style={draw=black!70, densely dashed, line width=0.45pt},
  act/.style={draw=black!70, fill=black!10, line width=0.45pt},
  msg/.style={fluxdonnee},
  rep/.style={fluxreponse},
  % Le refus prend l'aspect de la convention : trait epais a l'encre d'accent,
  % barre de la croix. Il etait en tirets, aspect reserve au flux asynchrone.
  ref/.style={fluxrefuse},
  et/.style={font=\scriptsize, fill=white, inner sep=1pt, align=center},
  eg/.style={font=\scriptsize, fill=white, inner sep=1pt, align=left},
  num/.style={circle, draw=black!70, fill=white, line width=0.4pt,
              inner sep=0.5pt, minimum size=3.4mm, font=\scriptsize},
  stn/.style={circle, draw=black!70, fill=black!10, line width=0.5pt,
              inner sep=0.5pt, minimum size=3.8mm, font=\scriptsize\bfseries},
  note/.style={draw=black!70, fill=black!4, align=left, inner sep=3.5pt,
               font=\scriptsize},
  lead/.style={draw=black!55, densely dotted, line width=0.4pt},
  cart/.style={draw=black!45, rounded corners=2pt, fill=black!2, align=left,
               inner sep=4pt, font=\scriptsize}]

% ===================== participants ==========================================
\node[part, text width=1.40cm] (nv) at (1.15,-0.55) {Navigateur};
\node[part, text width=2.45cm] (tb) at (3.70,-0.55) {Tableau de bord\\(passerelle applicative)};
\node[part, text width=1.90cm] (ap) at (6.30,-0.55) {API de supervision};
\node[part, text width=1.60cm] (bd) at (8.60,-0.55) {Base des comptes};

% ===================== lignes de vie =========================================
\draw[lif] (1.15,-1.05) -- (1.15,-14.20);
\draw[lif] (3.70,-1.05) -- (3.70,-14.20);
\draw[lif] (6.30,-1.05) -- (6.30,-14.20);
\draw[lif] (8.60,-1.05) -- (8.60,-14.20);

% ===================== bande d'etats, a droite ===============================
\draw[black!70, line width=0.5pt] (10.20,-1.05) rectangle (15.70,-10.55);
\draw[black!55, line width=0.4pt] (10.20,-4.70) -- (15.70,-4.70);
\draw[black!55, line width=0.4pt] (10.20,-6.40) -- (15.70,-6.40);
\draw[black!55, line width=0.4pt] (10.20,-8.15) -- (15.70,-8.15);
\draw[black!55, line width=0.4pt] (10.20,-9.25) -- (15.70,-9.25);
\node[anchor=north west, font=\scriptsize\itshape, text=black!55, inner sep=0pt]
      at (10.34,-1.16) {états de la session du navigateur};

\node[stn] at (10.47,-1.85) {1};
\node[anchor=north west, text width=4.70cm, align=left, inner sep=0pt]
      at (10.78,-1.70) {non authentifié};

\node[stn] at (10.47,-5.00) {2};
\node[anchor=north west, text width=4.70cm, align=left, inner sep=0pt]
      at (10.78,-4.85) {mot de passe vérifié, second facteur en attente ---
      \textbf{jeton intermédiaire, 5~min, sans revendication de rôle}};

\node[stn] at (10.47,-6.70) {3};
\node[anchor=north west, text width=4.70cm, align=left, inner sep=0pt]
      at (10.78,-6.55) {code vérifié, rôle relu en base};

\node[stn] at (10.47,-8.45) {4};
\node[anchor=north west, text width=4.70cm, align=left, inner sep=0pt]
      at (10.78,-8.30) {session active, 60~min, autorisation par rôle à
      chaque appel};

\node[stn] at (10.47,-9.55) {5};
\node[anchor=north west, text width=4.70cm, align=left, inner sep=0pt]
      at (10.78,-9.40) {expirée --- \emph{aucune révocation serveur~: la
      déconnexion n'invalide pas le jeton}};

% retour de l'etat 5 vers l'etat 1, dans le couloir a gauche de la bande
% Retour de l'etat « expiree » vers « non authentifie » : c'est une transition
% d'etat, non un refus — elle ne prend donc pas le trait de refus.
\draw[-{Latex[length=1.6mm]}, draw=menaltrait, line width=0.5pt]
     (10.20,-9.90) -- (9.95,-9.90) -- (9.95,-2.20) -- (10.18,-2.20);

% reperes d'etat portes par la ligne de vie du navigateur
\node[stn] at (0.62,-2.85)  {1}; \draw[lead] (0.81,-2.85)  -- (1.04,-2.85);
\node[stn] at (0.62,-6.15)  {2}; \draw[lead] (0.81,-6.15)  -- (1.04,-6.15);
\node[stn] at (0.62,-7.25)  {3}; \draw[lead] (0.81,-7.25)  -- (1.04,-7.25);
\node[stn] at (0.62,-8.70)  {4}; \draw[lead] (0.81,-8.70)  -- (1.15,-8.70);
\node[stn] at (0.62,-9.90) {5}; \draw[lead] (0.81,-9.90) -- (1.15,-9.90);

% ===================== barres d'activation ===================================
\draw[act] (1.04,-1.75) rectangle (1.26,-8.00);
\draw[act] (3.59,-1.75) rectangle (3.81,-4.30);
\draw[act] (6.19,-2.40) rectangle (6.41,-4.30);
\draw[act] (6.19,-5.25) rectangle (6.41,-8.00);
\draw[act] (8.49,-3.35) rectangle (8.71,-3.60);
\draw[act] (8.49,-6.05) rectangle (8.71,-6.30);
\draw[act] (8.49,-6.80) rectangle (8.71,-7.05);

% ===================== messages du parcours nominal ==========================
% 1
\draw[msg] (1.26,-1.80) -- (3.59,-1.80);
\node[num] at (1.62,-1.80) {1};
\node[et] at (2.65,-1.44) {adresse et mot de passe};

% 2
\draw[msg] (3.81,-2.45) -- (6.19,-2.45);
\node[num] at (4.15,-2.45) {2};
\node[et] at (5.20,-2.18) {relais};

\node[note, text width=1.95cm, anchor=north west] (n2) at (1.34,-2.78)
      {\textbf{aucun témoin de connexion posé à cette étape}};
\fill[white] (n2.north east) -- ++(-0.26,0) -- ++(0,-0.26) -- cycle;
\draw[black!70, line width=0.5pt]
      ([xshift=-0.26cm]n2.north east) -- ([yshift=-0.26cm]n2.north east);
\draw[lead] (n2.east) -- (4.95,-2.56);

% 3
\draw[msg] (6.41,-3.45) -- (8.49,-3.45);
\node[num] at (6.72,-3.45) {3};
\node[et] at (8.05,-2.98) {vérification du mot de passe\\(hachage adaptatif)};

% 4
\draw[rep] (6.19,-4.20) -- (3.81,-4.20);
\node[num] at (5.88,-4.20) {4};
\node[et] at (5.00,-4.58) {second facteur requis + \textbf{jeton intermédiaire}};

% 5
\draw[msg] (1.26,-5.30) -- (6.19,-5.30);
\node[num] at (1.62,-5.30) {5};
\node[et] at (4.00,-4.98) {jeton intermédiaire + code à six chiffres};

% 6
\draw[msg] (6.41,-6.15) -- (8.49,-6.15);
\node[num] at (6.72,-6.15) {6};
\node[et] at (7.95,-5.40) {déchiffrement du secret\\partagé, vérification\\(pas de 30~s,\\tolérance $\pm$1 pas)};

% 7
\draw[msg] (6.41,-6.90) -- (8.49,-6.90);
\node[num] at (6.72,-6.90) {7};
\node[et] at (7.45,-6.52) {\textbf{relecture du rôle}};

% 8
\draw[rep] (6.19,-7.90) -- (1.26,-7.90);
\node[num] at (5.88,-7.90) {8};
\node[et] at (3.80,-7.54) {jeton d'accès, 60~min,\\témoin inaccessible au script};

% ===================== fragment des refus ====================================
\draw[black!70, line width=0.5pt] (0.42,-10.95) rectangle (15.70,-14.20);
\fill[black!10] (0.42,-10.95) -- (1.92,-10.95) -- (1.92,-11.18)
                -- (1.74,-11.40) -- (0.42,-11.40) -- cycle;
\draw[black!70, line width=0.5pt] (0.42,-10.95) -- (1.92,-10.95) -- (1.92,-11.18)
                -- (1.74,-11.40) -- (0.42,-11.40) -- cycle;
\node[anchor=west, inner sep=0pt] at (0.55,-11.175) {\textbf{refus}};

\draw[act] (1.04,-11.75) rectangle (1.26,-13.95);
\draw[act] (6.19,-11.75) rectangle (6.41,-13.95);

% r1 -- mot de passe faux
\draw[ref] (6.19,-11.85) -- node[croixrefus]{} (1.26,-11.85);
\node[eg, text width=8.60cm, anchor=west] at (6.55,-11.85)
      {\textbf{mot de passe faux} --- 401 : réponse identique que le second
       facteur soit actif ou non, \textbf{aucun oracle}};
% r2 -- compte desactive
\draw[ref] (6.19,-12.55) -- node[croixrefus]{} (1.26,-12.55);
\node[eg, anchor=west] at (6.55,-12.55) {\textbf{compte désactivé} --- 403};
% r3 -- jeton d'acces presente au second facteur
\draw[ref] (6.19,-13.20) -- node[croixrefus]{} (1.26,-13.20);
\node[eg, anchor=west] at (6.55,-13.20)
      {\textbf{jeton d'accès présenté à la vérification du second facteur} --- 401};
% r4 -- jeton intermediaire presente a une route metier
\draw[ref] (6.19,-13.85) -- node[croixrefus]{} (1.26,-13.85);
\node[eg, anchor=west] at (6.55,-13.85)
      {\textbf{jeton intermédiaire présenté à une route métier} --- 401};

% ===================== notation ==============================================
\node[anchor=north west, inner sep=0pt, align=left, text width=15.28cm]
      at (0.42,-14.60)
  {\textbf{Notation}\;
   \tikz[baseline=-0.55ex]{\draw[msg] (0,0) -- (0.62,0);}\,appel synchrone \quad
   \tikz[baseline=-0.55ex]{\draw[rep] (0,0) -- (0.62,0);}\,réponse \quad
   \tikz[baseline=-0.55ex]{\draw[ref] (0,0) -- (0.62,0);}\,refus, marqué d'une croix \quad
   \tikz[baseline=-0.55ex]{\node[stn, minimum size=3.2mm]{\ };}\,état de la session};

\end{tikzpicture}
\caption[Authentification à deux facteurs]{Authentification à deux facteurs : cinq états de la
session et le jeton intermédiaire qui ne porte aucun rôle. Entre le mot de passe et le code à usage
unique, la session existe \textbf{authentifiée mais non autorisée} ; c'est la séparation des types
de jeton --- non le second facteur --- qui interdit d'utiliser ce demi-passage.}
\label{fig:seq-mfa}
\end{figure}

\section{Décisions d'architecture et écarts d'implémentation}

\subsection{Registre des décisions}

Le registre prouve qu'une architecture est délibérée : chaque décision y figure avec son statut
(\cref{tab:registre-decisions}), sa justification complète étant reportée en
annexe~\ref{ann:conception-detaillee}. Il en compte dix-sept, de D00 à D16 ; la première, le choix
d'un fournisseur unique, conditionne toutes les suivantes.

\begin{table}[htbp]
\caption[Registre synthétique des décisions d'architecture]{Registre synthétique des décisions d'architecture (justification complète en \cref{ann:conception-detaillee}).}
\label{tab:registre-decisions}
\centering
\footnotesize
\setlength{\tabcolsep}{3.5pt}
\renewcommand{\arraystretch}{1.05}
\begin{tabular}{|P{1.0cm}|P{10.5cm}|P{3.6cm}|}
\hline
\# & Décision & Statut \\
\hline
D00 & Google Cloud comme fournisseur unique du socle & Validée \\ \hline
D01 & Exécution de conteneurs sans serveur comme unique plateforme & Validée \\ \hline
D02 & Point d'entrée unique avec filtrage applicatif & Validée \\ \hline
D03 & Base de données managée, adresse privée, restauration à un instant donné & Validée \\ \hline
D04 & Entrepôt de données généraliste comme système de supervision & Validée \\ \hline
D05 & Modèle de représentation non génératif pour l'enrichissement & Validée \\ \hline
D06 & Recherche vectorielle déportée dans l'entrepôt & Validée \\ \hline
D07 & Registre d'images unique, retrait du second registre & Retouchée \\ \hline
D08 & Fédération d'identité pour la chaîne de livraison & Ajoutée \\ \hline
D09 & Détection de secrets en tête de chaîne & Ajoutée \\ \hline
D10 & Autorisation applicative par rôles conservée & Validée \\ \hline
D11 & Sortie réseau : route par défaut, aucune règle de sortie déclarée & Réserve, levée par D15 \\ \hline
D12 & Infrastructure en code comme standard unique de provisionnement & Validée, renforcée \\ \hline
D13 & Passage de la base de données en haute disponibilité régionale & Révision de D03
(08/08/2026) \\ \hline
D14 & Réécriture du tableau de bord sur cadre à rendu serveur avec passerelle applicative &
Validée \\ \hline
D15 & Sortie réseau directe, refus par défaut en sortie et pare-feu ciblé par identité &
\textbf{Appliquée} du 11 au 25/08/2026 (\cref{subsec:reseau-cible}) \\ \hline
D16 & Isolation des données proportionnée à la sensibilité : instance dédiée au locataire
porteur de donnée irréparable & Arrêtée le 29/08/2026, non appliquée (\cref{subsec:isolation-donnees}) \\
\hline
\end{tabular}
\end{table}

\textbf{La décision D13 illustre la fonction du registre : elle contredit D03}, contradiction
enregistrée telle quelle plutôt que masquée par une réécriture rétroactive. Un second registre,
tenu au fil de la réalisation, compte quinze entrées du 30/07/2026 au 19/08/2026 ; les deux sont
complémentaires (détail en annexe~\ref{ann:conception-detaillee}).

\subsection{Composants évalués puis écartés}
\label{tab:composants-ecartes}

Sept éléments ont été évalués puis écartés en application de PR1. Un seul l'a été pour un motif de
sécurité : le \textbf{modèle de langue génératif}, dont la sortie non déterministe et la surface
d'injection sont rédhibitoires pour une chaîne de preuve --- \textbf{ce refus ne sera pas révisé si
le contexte change}. Les six autres --- périmètre anti-exfiltration, orchestrateur et maillage de
services, cache distribué, déploiement multi-région, signature des images au déploiement, passerelle
d'authentification --- relèvent d'un arbitrage de coût ou d'une redondance avec un contrôle déjà
présent (\cref{ann:composants-ecartes}).

\subsection{Écarts entre la conception et l'implémentation}

Tout projet d'infrastructure produit des écarts entre l'architecture décrite et le système déployé.
La question n'est pas de savoir s'il y en a, mais s'ils sont connus, justifiés, datés lorsqu'ils
sont clos et mesurés dans leurs conséquences ; ceux encore ouverts portent, à la place d'une date,
leur état d'avancement, et ils figurent ici plutôt qu'en annexe
(\cref{tab:ecarts-implementation}). Deux d'entre eux, É10 et É11, portaient sur ce que les
vues laissaient croire : les vues ont été corrigées, les écarts subsistent parce qu'ils décrivent
une propriété du système et non un défaut de dessin.

\begin{table}[htbp]
\caption{Écarts entre conception et implémentation.}
\label{tab:ecarts-implementation}
\centering
\footnotesize
\setlength{\tabcolsep}{3.5pt}
\renewcommand{\arraystretch}{1.05}
\begin{tabular}{|P{0.75cm}|P{7.70cm}|P{1.95cm}|P{1.85cm}|P{2.34cm}|}
\hline
\# & Écart et cause & Date ou état & Exigence & Statut \\
\hline
É1 & Journalisation des refus du pare-feu et de la passerelle de sortie initialement absente &
19/08/2026 & EX20 & Corrigé, déployé et vérifié \\ \hline
É2 & Les deux sous-réseaux prévus, public \emph{et} privé, sont inutilisés : aucune ressource ne
les occupait, aucune machine virtuelle n'existait dans le dépôt et les sorties du module réseau
n'étaient consommées par aucun autre module ; les deux sous-réseaux ont été supprimés au titre de
D15 & 24/08/2026 & --- & Corrigé, appliqué et vérifié \\ \hline
É3 & Service d'encodage partage l'identité de la tâche de détection planifiée, faute d'identité
dédiée & 19/08/2026 & EX5, EX14 & Corrigé le 19/08/2026, vérifié \\ \hline
É4 & Mécanisme de connexion au réseau privé différent de celui prévu (confort d'un composant
managé) : le connecteur d'accès sans serveur a été supprimé et remplacé par la sortie réseau
directe au titre de D15 & 15/08/2026 & --- & Corrigé, appliqué et vérifié \\ \hline
É5 & Cadence de la tâche d'enrichissement plus lente que prévu (arbitrage coût/latence) & mise
en œuvre & --- & Assumé \\ \hline
É6 & Cache d'encodage instrumenté mais non utilisé (non terminé) & mise en œuvre & --- &
À terminer \\ \hline
É7 & Politiques d'organisation décrites en code mais inapplicables sans organisation GCP &
19/08/2026 & EX7 & Ouvert, contrôle compensatoire par IAM et IaC \\ \hline
É8 & Déploiement applicatif référencé par une étiquette égale au SHA du commit, et non par une
empreinte d'image stricte & 19/08/2026 & EX9 & Assumé, provenance à renforcer \\ \hline
É9 & Cloisonnement du tableau de bord non appliqué par les autorisations : l'interface passe par
l'API de supervision dans son code, mais son identité de service conserve un droit de lecture sur
le jeu de données de supervision et le droit d'exécuter des requêtes ; la séparation est
applicative et non structurelle & 27/08/2026 & EX5 & Ouvert \\ \hline
É10 & Zone de noms de domaine hors du périmètre décrit en code : aucune ressource de zone n'est
déclarée, les enregistrements sont créés à la main chez le bureau d'enregistrement, par exception
au principe de description en code (PR4) ; les vues la portent en dépendance externe, jamais en
composant de la couche de périmètre & 27/08/2026 & --- & Ouvert
\\ \hline
É11 & L'accès à l'entrepôt de supervision ne traverse pas le réseau privé : depuis la refonte, les
appels sortent bien par l'interface de la charge, mais la règle qui les admet (E4) désigne la plage
des interfaces du fournisseur et non l'entrepôt du socle --- l'autorisation reste portée par
l'identité seule ; la vue de déploiement distingue ce chemin par le style du lien & 27/08/2026 &
EX5 & Assumé, documenté \\
\hline
\end{tabular}
\end{table}

\textbf{L'écart É3 a été le plus significatif}, entorse au principe d'identité dans la couche même
qui l'illustre : l'identité dédiée a été déployée et vérifiée le 19/08/2026. La séparation est
établie par la description en code et par la revue, non par une épreuve dynamique --- le contrôle
d'usurpation qui l'établirait en propre ne figure pas parmi les six protocoles du plan de validation
continue, et la réserve est portée comme telle. La refonte réseau en a clos deux autres, É2 et É4 ;
les trois derniers, relevés le 27/08/2026, laissent É9 et É10 ouverts et É11 assumé.

\section*{Conclusion du chapitre}
\addcontentsline{toc}{section}{Conclusion du chapitre}

Quatre principes ont guidé l'architecture qui prend en charge les vingt exigences du \cref{ch:besoins-menaces} :
justification de chaque composant, identité comme plan de contrôle plutôt que comme couche,
intégrité de la chaîne de preuve, description exclusive de l'infrastructure en code. Le modèle
comporte sept niveaux L1--L7, cinq couches et deux plans transversaux, chacun décrit jusqu'à son
point de vigilance. Huit vues d'architecture et trois diagrammes de comportement rendent ces choix
vérifiables plutôt que déclaratifs --- dont la topologie réseau, qui répond à la question de la
sortie, et le modèle entité-association, qui \emph{montre} pourquoi le rattachement sémantique ne
rejoint pas le calcul de score --- et le \cref{tab:recette-vues} énonce, pour les six vues qui
portent une décision opposable, le point de vue servi, les critères de recette datés avec leur
protocole, et la réserve. Sept flux
structurent la vue dynamique ; le sixième, hypothèse assortie de son protocole de réfutation,
relie les vulnérabilités détectées à la livraison aux techniques d'attaque observées en recette.
Dix-sept décisions ont été enregistrées, de D00 à D16, dont une révision explicite d'une décision
antérieure, et sept composants écartés. Les onze écarts recensés sont datés lorsqu'ils sont clos
et rattachés aux exigences qu'ils affectent ; quatre sont clos et vérifiés, deux le 19/08/2026 et
deux par la refonte réseau, les 15 et 24/08/2026.

% -- Restauration des valeurs par défaut de placement et d'espacement des flottants.
\renewcommand{\topfraction}{0.7}
\renewcommand{\bottomfraction}{0.3}
\renewcommand{\textfraction}{0.2}
\renewcommand{\floatpagefraction}{0.5}
\setcounter{topnumber}{2}
\setcounter{totalnumber}{3}
\setlength{\textfloatsep}{20pt plus 2pt minus 4pt}
\setlength{\floatsep}{12pt plus 2pt minus 2pt}
\setlength{\intextsep}{12pt plus 2pt minus 2pt}
\setlength{\abovecaptionskip}{10pt}
\setlength{\belowcaptionskip}{0pt}
